

\section{Povsem pozitivne preslikave}\label{sec:2}

V tem magistrskem delu bomo, če ni izrecno napisano drugače,
za vse $C^*$-algebre predpostavili, da so enotske in nad kompleksnimi skalarji.
Prav tako bomo vsi $*$-homomorfizmi $C^*$-algeber $\varphi: \mathfrak{A} \to \mathfrak{B}$ enotski,
torej bodo slikali $\varphi(1) = 1$. Vse vektorske prostore bomo obravnavali nad kompleksnimi skalarji. 

Naj bo $\mathfrak{A}$ $C^*$-algebra. Definiramo $M_n (\mathfrak{A})$ kot algebro $n \times n$ matrik z elementi v $\mathfrak{A}$, z običajnim matričnim seštevanjem in množenjem.
Na $M_n(\mathfrak{A})$ lahko definiramo matrično involucijo
$$[a_{ij}]_{i, j} ^* := [a_{ji} ^*]_{i, j},$$
s katero $M_n (\mathfrak{A})$ postane $*$-algebra.

\begin{definicija}
    Naj bosta $\mathfrak{A}, \mathfrak{B}$ $C^*$-algebri in $\phi: \mathfrak{A} \to \mathfrak{B}$ linearna preslikava.
    Za vsak $n \in \N$ definiramo 
    $$\phi^{(n)}: M_n (\mathfrak{A}) \to M_n (\mathfrak{B}),\quad [a_{ij}]_{i, j} \mapsto [\phi(a_{ij})]_{i, j}.$$
    Preslikavi $\phi^{(n)}$ pravimo $n$-ta \emph{ampliacija} $\phi$.
\end{definicija}

Če je $\mathfrak{A}$ $C^*$-algebra, potem po Gelfand--Naimark--Segalovem izreku (\cite[Theorem VIII.5.14., str. 250]{conway1990}, v nadaljevanju: izrek GNS)
obstaja Hilbertov prostor $\mathcal{H}$ in injektivna upodobitev $\pi: \mathfrak{A} \to \bh$, kar inducira injektivni $*$-homomorfizem 
$$\pi^{(n)}: M_n (\mathfrak{A}) \to M_n (\bh)$$
za vsak $n \in \N$. To pomeni, da lahko $M_n (\mathfrak{A})$ identificiramo z
$*$-podalgebro $C^*$-algebre $M_n (\bh)$.
Ker je slika $\pi^{(n)} (M_n (\mathfrak{A}))$ zaprta v $M_n (\bh)$,
je tudi $C^*$-podalgebra v $M_n (\bh)$.
To nam inducira normo na $M_n (\mathfrak{A})$, s katero $M_n (\mathfrak{A})$ postane $C^*$-algebra.
Spomnimo, da vsaka $*$-algebra premore največ eno normo, za katero postane $C^*$-algebra (glej recimo \cite[Corollary 2.1.2., str. 37]{murphy1990}),
kar pomeni, da je zgoraj opisana norma na $M_n (\mathfrak{A})$ kanonična oziroma neodvisna od izbora upodobitve $\pi$.
Kot bomo videli, nam kanonične norme na $M_n (\mathfrak{A})$ povedo dodatne informacije o začetni $C^*$-algebri $\mathfrak{A}$.

\begin{definicija}
    Element $a \in \mathfrak{A}$ je \emph{pozitiven}, če obstaja $b \in \mathfrak{A}$, tako da je $a = b^* b$.
    Množico vseh pozitivnih elementov v $\mathfrak{A}$ označimo z $\mathfrak{A}_+$.
\end{definicija}

\begin{opomba}
    Množica $\mathfrak{A}_+$ je \emph{konveksen stožec}: to pomeni, da je $\mathfrak{A}_+ + \mathfrak{A}_+ \subseteq \mathfrak{A}_+$,
    $c \mathfrak{A}_+ \subseteq \mathfrak{A}_+$ za $\alpha \geq 0$ ter $\mathfrak{A}_+ \cap - \mathfrak{A}_+ = \{0\}$.
\end{opomba}

\begin{definicija}
    Naj bosta $\mathfrak{A} , \mathfrak{B}$ $C^*$-algebri in $\varphi: \mathfrak{A} \to \mathfrak{B}$ linearna preslikava.
    \begin{enumerate}
        \item $\varphi$ je \emph{pozitivna}, če je $\varphi(\mathfrak{A}_+) \subseteq \mathfrak{B}_+$;
        \item $\varphi$ je \emph{$n$-pozitivna}, če je $\varphi^{(n)}$ pozitivna;
        \item $\varphi$ je \emph{povsem pozitivna}, če je $n$-pozitivna za vsak $n \in \N$.
    \end{enumerate}
    Množico vseh (enotskih) povsem pozitivnih preslikav označimo s $\PP (\mathfrak{A}, \mathfrak{B})$ (oziroma $\EPP (\mathfrak{A} , \mathfrak{B})$).
\end{definicija}

\begin{opomba}
    \begin{itemize}
        \item Če je $\varphi$ $n$-pozitivna, potem je tudi $k$-pozitivna za vse $k \leq n$.
        \item Kompozitum pozitivnih ($n$-pozitivnih, povsem pozitivnih) preslikav je prav tako pozitiven ($n$-pozitiven, povsem pozitiven).
        \item Če sta $\varphi, \psi: \mathfrak{A} \to \mathfrak{B}$ pozitivni ($n$-pozitivni, povsem pozitivni) preslikavi in $\alpha, \beta \geq 0$,
        potem je tudi $\alpha \varphi + \beta \psi$ pozitivna ($n$-pozitivna, povsem pozitivna).
    \end{itemize}
\end{opomba}

V nadaljevanju si bomo oglejmo nekaj primerov povsem pozitivnih preslikav
in njihovih osnovnih lastnosti. Model povsem pozitivnih preslikav so pozitivni funkcionali. 

\begin{zgled}\label{zgl:1.3}
    Dokažimo, da je vsak pozitivni linearni funkcional $\varphi$ na $C^*$-algebri $\mathfrak{A}$ povsem pozitiven.
    Dovolj je pokazati, da je $$\varphi^{(n)} : M_n (\mathfrak{A}) \to M_n (\C)$$
    pozitivna preslikava. Ključen je uvid, da za popolnoma poljuben funkcional $\varphi$, matriko $A = [a_{ij}]_{i, j} \in M_n (\mathfrak{A})$
    in vektorja $$\alpha = \begin{bmatrix}
        \alpha_1 \\ \vdots \\ \alpha_n
    \end{bmatrix},\quad \beta = \begin{bmatrix}
        \beta_1 \\ \vdots \\ \beta_n
    \end{bmatrix} \in \C^n$$ velja 
    \begin{align*}
        \beta^* \varphi^{(n)} (A) \alpha &= \begin{bmatrix}
            \beta_1  \\ \vdots \\ \beta_n 
        \end{bmatrix}^*
        \begin{bmatrix}
            \varphi (a_{11}) & \cdots & \varphi(a_{1n})\\
            \vdots & & \vdots\\
            \varphi (a_{n1}) & \cdots & \varphi(a_{nn})
        \end{bmatrix}
        \begin{bmatrix}
            \alpha_1 \\ \vdots \\ \alpha_n
        \end{bmatrix} \\
        &= \sum_{i, j} \overline{\beta_i} \varphi(a_{ij}) \alpha_j  = \varphi \left(\sum_{i, j} \overline{\beta_i} a_{ij} \alpha_j \right)\\
        &= \varphi(\beta^* A \alpha).
    \end{align*}
    Če je sedaj matrika $A \in M_n (\mathfrak{A})$ pozitivna, potem je za vsak $\alpha \in \C^n$
    tudi $\alpha^* A \alpha \in \mathfrak{A}$ pozitiven element in velja 
    $$\alpha^* \varphi^{(n)}(A) \alpha = \varphi(\alpha^* A \alpha) \geq 0.$$
    To pomeni, da je $\varphi^{(n)} (A) \in M_n (\C)$ pozitivna matrika, zato je $\varphi^{(n)}: M_n (\mathfrak{A}) \to M_n (\C)$ pozitivna preslikava.
\end{zgled}

To nam omogoča, da na povsem pozitivne preslikave gledamo kot na posplošitve stanj,
pri čemer je kodomena poljubna $C^*$-algebra. Izkazalo se bo, da lahko teorijo stanj posplošimo na 
povsem pozitivne preslikave.

\begin{zgled}
    Če je $\varphi: \mathfrak{A} \to \mathfrak{B}$ $*$-homomorfizem, je $\varphi$ pozitiven, saj 
    $\varphi(a^* a) = \varphi(a) ^* \varphi(a)$. Še več, v tem primeru je tudi $\varphi^{(n)}$ $*$-homomorfizem,
    zato je $\varphi^{(n)}$ pozitivna za vsak $n\in \N$. Od tod sledi,
    da je vsak $*$-homomorfizem med dvema $C^*$-algebrama povsem pozitiven.
\end{zgled}

\begin{zgled}\label{zgl:1.4}
    Obstajajo pozitivne preslikave, ki niso povsem pozitivne.
    Primer take preslikave je $$\varphi: M_2 (\C) \to M_2 (\C),\quad A \mapsto A^\top$$
    (tukaj operacija $A^\top$ pomeni realno transpozicijo matrike brez konjugiranja, saj mora biti $\varphi$ linearna preslikava).
    Očitno je $\varphi$ pozitivna preslikava, kar pa ne velja za njeno drugo ampliacijo
    $$\varphi^{(2)}:  M_4 (\C) \to M_4 (\C),\quad \begin{bmatrix}
        A_{11} & A_{12}\\
        A_{21} & A_{22}
    \end{bmatrix} \mapsto \begin{bmatrix}
        A_{11}^\top & A_{12}^\top\\
        A_{21}^\top & A_{22}^\top
    \end{bmatrix}.$$ Ta namreč slika pozitivno matriko
    $$\begin{bmatrix}
        1 & 0 & 0 & 1\\
        0 & 0 & 0 & 0\\
        0 & 0 & 0 & 0\\
        1 & 0 & 0 & 1
    \end{bmatrix} = \begin{bmatrix}
        1 \\ 0 \\ 0 \\ 0
    \end{bmatrix} \begin{bmatrix}
        1 \\ 0 \\ 0 \\ 0
    \end{bmatrix}^* \geq 0$$
    v
    $$\begin{bmatrix}
        1 & 0 & 0 & 0\\
        0 & 0 & 1 & 0\\
        0 & 1 & 0 & 0\\
        0 & 0 & 0 & 1
    \end{bmatrix},$$
    ki ima negativno lastno vrednost $-1$ (in zato seveda ne more biti pozitivna).
\end{zgled}

\begin{zgled}\label{zgl:1.5}
    Naj bosta $\mathcal{H}, \mathcal{K}$ Hilbertova prostora in $V \in \mathcal{B}(\mathcal{K}, \mathcal{H})$
    omejen operator. Dokazujemo, da je preslikava
    $$\varphi: \bh \to  \bk,\quad A \mapsto V^* A V$$
    povsem pozitivna. Vzemimo pozitivno matriko $[A_{ij}]_{i, j} \in M_n (\bh)_+$.
    Potem je 
    \begin{align*}
        \varphi^{(n)} ([A_{ij}]_{i, j}) &= \begin{bmatrix}
            \varphi(A_{11}) & \cdots & \varphi(A_{1n})\\
            \vdots & & \vdots\\
            \varphi(A_{n1}) & \cdots & \varphi(A_{nn})\\
        \end{bmatrix}\\
        &= \begin{bmatrix}
            V^* A_{11} V & \cdots & V^* A_{1n} V\\
            \vdots & & \vdots\\
            V^* A_{n1} V & \cdots & V^* A_{nn} V\\
        \end{bmatrix}\\
        &= \begin{bmatrix}
            V & &\\
             & \ddots & \\
             & & V
        \end{bmatrix}^* \begin{bmatrix}
            A_{11} & \cdots & A_{1n}\\
            \vdots & & \vdots\\
            A_{n1} & \cdots & A_{nn}\\
        \end{bmatrix} \begin{bmatrix}
            V & &\\
             & \ddots & \\
             & & V
        \end{bmatrix}\\
    \end{align*}
    pozitivna matrika, torej je $\varphi ^{(n)}$ pozitivna preslikava. 
    Posledično je $\varphi$ povsem pozitivna.

    Če je $\mathfrak{A}$ $C^*$-algebra in $\psi: \mathfrak{A} \to \bh$ povsem pozitivna preslikava,
    potem je tudi kompozicija 
    $$
        \mathfrak{A} \to  \bk,\quad A \mapsto V^* \psi (a) V 
    $$
    povsem pozitivna.
\end{zgled}

Razdelek zaključimo z uporabno karakterizacijo povsem pozitivnih preslikav \eqref{eq:intro1}
iz uvoda.


\begin{comment}
\begin{zgled}
    Naj bo $\psi: \mathfrak{A} \to \mathcal{B}$ povsem pozitivna in $b \in B$.
    Definirajmo $$\varphi :A \to B ,\quad a \mapsto b^* \psi(a) b.$$
    Dokazujemo, da je $\varphi$ povsem pozitivna. Za vsak $n \in \N$ in $[a_{ij}]_{i, j} \in M_n (A)_+$
    je namreč
    \begin{align*}
        \varphi^{(n)} ([a_{ij}]_{i, j}) &= [\varphi(a_{ij})]_{i, j} = \begin{bmatrix}
            b^* \psi (a_{11})  b & \cdots & b^* \psi (a_{1n}) b\\
            \vdots & \ddots & \vdots\\
            b^* \psi (a_{n1})  b & \cdots & b^* \psi (a_{nn}) b
        \end{bmatrix}\\
        &= \begin{bmatrix}
            b & & \\
             & \ddots & \\
              & & b
        \end{bmatrix}^* (\psi^{(n)} [(a_{ij})]_{i, j}) \begin{bmatrix}
            b & & \\
              & \ddots & \\
              & & b
        \end{bmatrix} \geq 0.
    \end{align*}
\end{zgled}
\end{comment}


\begin{lema}\label{lem:1.2}
    Če je $\mathfrak{A}$ $C^*$-algebra, potem je vsak pozitiven element $M_n (\mathfrak{A})$
    vsota $n$ pozitivnih matrik oblike $[a_i ^* a_j]_{i, j}$ za 
    neke $\{a_1, \dots, a_n\} \subseteq A$.
\end{lema}

\begin{dokaz}
    Najprej opazimo, da če imamo matriko $$R = \begin{bmatrix}
        0 & \cdots & 0\\
        \vdots & & \vdots\\
        a_1 & \cdots & a_n\\
        \vdots & & \vdots\\
        0 & \cdots & 0\\       
    \end{bmatrix} \in M_n (\mathfrak{A}),$$
    potem je $R^* R = [a_i^* a_j]_{i, j}$.
    Naj bo sedaj $P \in M_n (\mathfrak{A})_+$ poljubna pozitivna matrika, torej $P = B^* B$ za nek $B \in M_n (\mathfrak{A})$.
    Zapišemo $B = R_1 + \cdots + R_n$, kjer je $R_k$ matrika,
    ki ima $k$-to vrstico enako kot $B$, drugje pa je ničelna.
    Ker velja $R_k^* R_l = 0$ za poljubna različna indeksa $k \neq l$ iz $\{1, \dots, n\}$,
    sledi
    \begin{align*}
        P = B^* B &= (R_1 ^* + \cdots + R_n ^*) (R_1 + \cdots + R_n) \\
        &= R_1 ^* R_1 + \cdots + R_n ^* R_n. \qedhere
    \end{align*}
\end{dokaz}



\begin{lema}\label{lem:1.3}
    Naj bo $\mathfrak{B}$ $C^*$-algebra in $A = [a_{ij}]_{i, j} \in M_n(\mathfrak{B})$.
    Potem je $A$ pozitivna natanko tedaj, ko velja $$\begin{bmatrix}
        b_1 \\ \vdots \\ b_n
    \end{bmatrix}^* A \begin{bmatrix}
        b_1 \\ \vdots \\ b_n
    \end{bmatrix} = \sum_{i, j} ^n b_i ^* a_{ij} b_j \geq 0$$
    za poljubno $n$-terico $b_1, \dots, b_n \in \mathfrak{B}$.
\end{lema}

\begin{dokaz}
    Implikacija v levo je trivialna, zato dokazujemo implikacijo v desno.
    \begin{enumerate}
        \item  V prvem koraku dokažemo, da če je $A = [a_{ij}]_{i, j} \in M_n (\mathfrak{B})$ takšna matrika, da velja $$\begin{bmatrix}
        b_1 \\ \vdots \\ b_n
    \end{bmatrix}^* A \begin{bmatrix}
        b_1 \\ \vdots \\ b_n
    \end{bmatrix} = \sum_{i, j} ^n b_i ^* a_{ij} b_j = 0$$
    za poljubno $n$-terico $b_1, \dots, b_n \in \mathfrak{B}$, potem je $A = 0$.
    Dokazali bomo, da je vsaka komponenta $A$ ničelna. Če v zgornjo enačbo vstavimo $b_i = 1$ in $b_j = 0$ za vse $j \neq i$, potem dobimo
    $a_{ii} = 0$. Od tod sledi, da so vsi diagonalni elementi matrike $A$ ničelni. 
    Preostane še pokazati, da je $a_{ij} = 0$ za poljubna indeksa $i \neq j$.
    Če vstavimo $b_i := 1, b_j := c \in \mathfrak{B}$ in $b_k = 0$ za $k \neq i, j$, dobimo enakost
    $$a_{ii} + a_{ij} c + c^* a_{ji}  + c^* a_{jj} c = 0.$$
    Ker so diagonalni elementi ničelni, sledi
    $$a_{ij} c + c^* a_{ji} = 0.$$
    Če posebej vstavimo $c = 1$ in $c = i$, dobimo
    $$a_{ij} + a_{ji} = 0, \quad i a_{ij} - i a_{ji} = 0,$$
    od koder potem sledi $a_{ij} = a_{ji} = 0$. S tem smo dokazali, da je $A = 0$.
        \item V drugem koraku dokažemo, da če je $$\begin{bmatrix}
            b_1 \\ \vdots \\ b_n
        \end{bmatrix}^* A \begin{bmatrix}
            b_1 \\ \vdots \\ b_n
        \end{bmatrix} = \sum_{i, j} ^n b_i ^* a_{ij} b_j \in \mathfrak{B}$$
        sebiadjungiran element za vsako $n$-terico $b_1, \dots, b_n \in \mathfrak{B}$, potem je $A$ sebiadjungirana matrika.
        Opazimo namreč, da velja
        \begin{align*}
            \sum_{i, j} ^n b_i ^* a_{ij} b_i &= \left(\sum_{i, j} ^n b_i ^* a_{ij} b_i\right)^* \\
            &= \sum_{i, j} ^n b_j ^* a_{ij}^* b_i\\
            &=  \sum_{i, j} ^n b_i ^* a_{ji}^* b_j,
        \end{align*}
        od koder sledi
        $$\begin{bmatrix}
        b_1 \\ \vdots \\ b_n
    \end{bmatrix}^* (A - A^*) \begin{bmatrix}
        b_1 \\ \vdots \\ b_n
    \end{bmatrix} = \sum_{i, j} ^n b_i ^* (a_{ij} - a_{ji}^*) b_i = 0.$$ Po prvemu koraku tako dobimo $A - A^* = 0$.
        \item Nazadnje dokažemo, da če je $$\begin{bmatrix}
            b_1 \\ \vdots \\ b_n
        \end{bmatrix}^* A \begin{bmatrix}
            b_1 \\ \vdots \\ b_n
        \end{bmatrix} = \sum_{i, j} ^n b_i ^* a_{ij} b_j$$
        pozitiven za poljubno $n$-terico $b_1, \dots, b_n \in \mathfrak{B}$, potem je $A$ pozitivna matrika.
        Ker je $A$ po prejšnjem koraku sebiadjungirana, imamo po zveznem funkcijskem računu (glej recimo \cite[Section VIII.2, str. 236-240]{conway1990}) 
        dekompozicijo $A = A_+ - A_-$,
        kjer sta $A_+, A_-$ pozitivni matriki, tako da velja $A_+ A_- = A_- A_+ = 0$.
        Po predpostavki imamo
        $$\begin{bmatrix}
            b_1 \\ \vdots \\ b_n
        \end{bmatrix}^* A_- A A_- \begin{bmatrix}
            b_1 \\ \vdots \\ b_n
        \end{bmatrix}\geq 0,$$
        hkrati pa je izraz na levi strani enak
        $$- \begin{bmatrix}
            b_1 \\ \vdots \\ b_n
        \end{bmatrix}^* A_- ^3 \begin{bmatrix}
            b_1 \\ \vdots \\ b_n
        \end{bmatrix}\leq 0.$$ Od tod sledi
        $$- \begin{bmatrix}
            b_1 \\ \vdots \\ b_n
        \end{bmatrix}^* A_- ^3 \begin{bmatrix}
            b_1 \\ \vdots \\ b_n
        \end{bmatrix} = 0$$ za poljubno $n$-terico $b_1, \dots, b_n \in \mathfrak{B}$.
        Po prvem koraku nato dobimo $A_- ^3 = 0$, od koder po zveznem funkcijskem računu sledi $A_- = 0$.
        To pa pomeni, da je $A = A_+$ pozitivna matrika. \qedhere
    \end{enumerate}
\end{dokaz}

Iz lem \ref{lem:1.2} in \ref{lem:1.3} sledi naslednja trditev.

\begin{trditev}
    Preslikava $\varphi: \mathfrak{A} \to \mathfrak{B}$ je povsem pozitivna natanko tedaj,
    ko za vsak $n \in \N$ in $n$-terici $a_1, \dots, a_n \in \mathfrak{A}$ ter $b_1, \dots, b_n \in \mathfrak{B}$
    velja $$\sum_{i, j} ^n b_i ^* \varphi (a_i ^* a_j) b_j \geq 0.$$
\end{trditev}

\section{Operatorski sistemi}\label{sec:3}

V zgornjih zgledih smo obravnavali povsem pozitivne preslikave na $C^*$-algebrah.
V praksi pa imamo pogosto opravka s povsem pozitivnimi preslikavami, ki so definirane zgolj na 
podprostorih $C^*$-algebre, ki jim rečemo operatorski sistemi. V tem poglavju obravnavamo nekatere lastnosti 
pozitivnih in povsem pozitivnih preslikav na operatorskih sistemih. 

\begin{definicija}
    Naj bo $\mathfrak{A}$ enotska $C^*$-algebra. Njenemu podprostoru $\mathcal{M} \subseteq \mathfrak{A}$,
    ki vsebuje enico, pravimo \emph{operatorski prostor}. Operatorskemu prostoru $\mathcal{S} \subseteq \mathfrak{A}$,
    za katerega velja $\mathcal{S} = \mathcal{S}^*$, pravimo \emph{operatorski sistem}.
\end{definicija}

\begin{opomba}
    V splošni literaturi (glej na primer \cite{paulsen2002}) so operatorski sistemi običajno definirani zgolj kot podprostori $C^*$-algebre,
    torej brez zahteve, da vsebujejo enico. V tem magistrskem delu bomo obravnavali zgolj enotske operatorske prostore,
    zato dodamo to kot del definicije.
\end{opomba}

\begin{zgled}
    Če je $\mathcal{M}$ operatorski prostor, potem je $\mathcal{S} = \mathcal{M} + \mathcal{M}^*$ operatorski sistem.
\end{zgled}

Če je $\mathcal{M}$ (oziroma $\mathcal{S}$) operatorski prostor (sistem) v $\mathfrak{A}$,
potem je $M_n (\mathcal{M})$ operatorski prostor (sistem) v $M_n (\mathfrak{A})$ za vse $n \in \N$.
Pozitivni elementi v operatorskem prostoru so $\mathcal{M}_+ = \mathcal{M} \cap \mathfrak{A}_+$.
Na enak način kot za $C^*$-algebre podamo tudi definicijo pozitivne, $n$-pozitivne ter povsem pozitivne preslikave
$\varphi: \mathcal{M} \to \mathfrak{B}$ na operatorskem prostoru $\mathcal{M}$.

\begin{zgled}
    Če je $\varphi: \mathfrak{A} \to \mathfrak{B}$ povsem pozitivna preslikava
    in $\mathcal{M} \subseteq \mathfrak{A}$ poljuben operatorski prostor,
    potem je tudi $\varphi|_\mathcal{S}: \mathcal{S} \to \mathfrak{B}$ povsem pozitivna preslikava.
\end{zgled}

\begin{zgled}
    Če je $\mathcal{S} \subseteq \mathfrak{A}$ operatorski sistem in $\varphi: \mathcal{S} \to \C$
    pozitiven funkcional, potem je $\varphi$ povsem pozitivna preslikava. Dokaz poteka popolnoma enako kot v zgledu
    \ref{zgl:1.3}.
\end{zgled}

Operatorski sistemi imajo pomembno vlogo v teoriji se v teoriji povsem pozitivnih preslikav, saj je vsak njihov
element linearna kombinacija pozitivnih elementov.
Res, vsak element v operatorskemu sistemu $a \in \mathcal{S}$ je linearna kombinacija dveh sebiadjungiranih elementov 
$$a = \underbrace{\frac{a + a^*}{2}}_{\real a} + i \cdot \underbrace{\frac{a - a^*}{2i}}_{\imag a},$$
ki prav tako ležita v $\mathcal{S}$.
Poleg tega lahko vsak sebiadjungiran element zapišemo kot razliko dveh pozitivnih elementov 
$$a = \frac{1}{2} (\| a\| + a) - \frac{1}{2} (\| a\| - a),$$
ki ležita v $\mathcal{S}$. 

Operatorski sistemi so zato `naraven' objekt, na katerem so morfizmi povsem pozitivne preslikave.
Z oznako $\osys$ označimo kategorijo operatorskih sistemov, v kateri so morfizmi povsem pozitivne preslikave,
v $\osys_1$ pa enotske povsem pozitivne preslikave.

\begin{zgled}
    Za razliko od operatorskih sistemov lahko operatorski prostori vsebujejo le trivialne sebiadjungirane elemente, torej $\R \cdot 1$.
    Oglejmo si operatorski prostor $\mathcal{M} := \C [z] \subseteq C(\T)$.
    Najprej dokažimo, da je $\{z^k\}_{k \in \Z} \subseteq C(\T)$ linearno neodvisna množica.
    Denimo, da obstajajo cela števila $k_1 \leq k_2 \leq \cdots \leq k_l$ in skalarji $\alpha_1, \alpha_2, \dots, \alpha_l \in \C$,
    da velja $\sum_{j = 1} ^l \alpha_j z^{k_j} = 0$ na $\T$. Potem je tudi polinom 
    $\sum_{j = 1} ^l \alpha_j z^{k_j - k_1}$ ničelen na $\T$. Po osnovnem izreku algebre 
    ima vsak neničelen kompleksen polinom končno število ničel, zato mora biti $$\alpha_1 = \alpha_2 = \cdots = \alpha_l = 0.$$
    Sedaj lahko dokažemo, da so edini sebiadjungirani elementi v $\C[z]$ trivialni.
    Vzemimo poljuben sebiadjungiran polinom $p = \sum_{j = 0} ^l \alpha_j z^j$. Potem je
    \begin{align*}
        \sum_{j = 0} ^l \alpha_j z^j = \overline{\sum_{j = 0} ^l \alpha_j z^j} = \sum_{j = 0} ^l \overline{\alpha_j} \overline{z}^j = \sum_{j = 0} ^l \overline{\alpha_j} z^{-j}.
    \end{align*} 
    Od tod sledi $\alpha_0 \in \R$ in $\alpha_j = 0$ za $j \geq 1$.
\end{zgled}

Kot nam pove zgornji zgled, imajo operatorski sistemi lahko premalo 
pozitivnih elementov, da bi povsem pozitivne preslikave na njih bile zanimive.
Kljub temu pa obstaja vrsta preslikav
na operatorskem prostoru $\mathcal{M}$, ki jih lahko razširimo do povsem pozitivnih preslikav na operatorskem sistemu $\mathcal{M} + \mathcal{M}^*$.
Takšne preslikave bodo potem naravni morfizem na operatorskih prostorih.
V nadaljevanju tega razdelka bomo opisali to konstrukcijo.

\begin{lema}
    Naj bo $\mathcal{S}$ operatorski sistem in $\mathfrak{B}$ $C^*$-algebra.
    Če je $\varphi: \mathfrak{S} \to \mathfrak{B}$ pozitivna preslikava, potem je \emph{$*$-linearna}, torej
    $$\varphi(a^*) = \varphi(a) ^*,\quad \forall a \in \mathcal{S}.$$
\end{lema}

\begin{dokaz}
    Ta trditev očitno velja za $a \in \mathcal{S}_+$.
    Sedaj predpostavimo, da je $a \in \mathcal{S}$ sebiadjungiran (torej $a = a^*$). 
    Vemo, da lahko vsak tak $a$ zapišemo kot $a = \frac{1}{2} (\| a\| + a) - \frac{1}{2}(\| a\| - a)$.
    Sedaj dobimo
    \begin{align*}
        \varphi (a)^* &= \varphi \left(\frac{1}{2} \left(\| a \| + a \right) - \frac{1}{2} \left(\| a \| - a \right)\right)^*\\
        &= \varphi \left(\frac{1}{2} (\| a \| + a)\right)^* - \varphi \left(\frac{1}{2} \| a\| - a\right)^*\\
        &= \varphi \left(\frac{1}{2} (\| a \| + a)\right) - \varphi \left(\frac{1}{2} \| a\| - a \right)\\
        &= \varphi \left(\frac{1}{2} \left(\| a \| + a \right) - \frac{1}{2} \left(\| a \| - a \right)\right)\\
        &= \varphi(a) = \varphi(a^*).
    \end{align*}
    Tako za povsem splošen $a \in \mathcal{S}$ dobimo
    \begin{align*}
        \varphi(a)^* &= \varphi(\real a + i \cdot \imag a)^*\\
        &= \varphi(\real a)^* - i \varphi(\imag a)^*\\
        &= \varphi(\real a) - i \varphi (\imag a)\\
        &= \varphi(\real a - i \cdot \imag a) = \varphi(a^*). \qedhere
    \end{align*}
\end{dokaz}

\begin{lema}\label{lem:1.1}
    Naj bo $\mathcal{S}$ operatorski sistem in $\varphi: \mathcal{S} \to \C$ funkcional.
    Potem je $\varphi$ pozitiven natanko tedaj, ko je $\| \varphi\| = \varphi(1)$.
\end{lema}

\begin{dokaz}
    Ker za vsak $a \in \mathcal{S}_+$ velja
    $$-\| a\| \leq a\leq \| a\|,$$
    sledi 
    $$\varphi(a) \leq \varphi(\|a\|),\quad \forall a \in \mathcal{S}.$$
    Za $a, b \in \mathcal{S}$ definiramo seskvilinearno formo $\langle a, b\rangle := \varphi(b^* a)$,
    od koder po Cauchy-Schwarzevi neenakosti sledi
    \begin{equation}
        |\varphi(b^* a)|^2 \leq \varphi (b^* b) \cdot \varphi(a^* a). \label{eq:1.2}
    \end{equation}
    Če vstavimo $b = 1$, dobimo 
    \begin{align*}
        |\varphi(a)|^2 &\leq \varphi(a^* a) \cdot \varphi(1) \leq \varphi(\| a^* a\|) \cdot \varphi(1)\\
        & = \| a^* a\| \cdot \varphi(1)^2 = \| a\|^2 \cdot \varphi(1)^2.
    \end{align*}
    Če enačbo korenimo, dobimo $\| \varphi\| \leq \varphi(1)$.
    Ker je $\varphi(1) \leq \| \varphi\|$, smo implikacijo v desno dokazali.
    Za implikacijo v levo vzamemo
    $a \in \mathcal{S}_+$ in $\varphi (a) := \alpha + i \beta.$
    Za vsak $t \in \R$ je
  \begin{align*}
    \alpha^2 + (\beta + t \|\varphi\|)^2 &= |\alpha + i (\beta + t \varphi(1))|^2\\
    &= |\varphi (a + it)|^2\\
    &\leq \| a + it\|^2 \cdot \|\varphi\|^2 \\
    &= \left( \|a\|^2 + t^2 \right) \|\varphi\|^2. 
  \end{align*}
    Od tod sledi $2 \beta t \|\varphi\| \leq \|a\|^2 \cdot \|\varphi\|^2$.
    Ker je bil $t \in \R$ poljuben, imamo $\beta = 0$ in $\varphi(a) = \alpha \in \R$.
    Tako dobimo
    \begin{align*}
        \| a\| \cdot \|\varphi\| - \varphi(a) &= \varphi(\|a\| - a)\\
        &\leq \|\|a\| - a\| \cdot \|\varphi\|\\
        &\leq \| a\| \cdot \| \varphi \|,
    \end{align*}
    torej $\varphi(a) \geq 0$.
\end{dokaz}

Za poljubno pozitivno preslikavo $\varphi: \mathcal{S} \to \mathfrak{B}$ imamo na voljo slabšo oceno.

\begin{trditev}\label{trd:1.1}
    Naj bo $\mathcal{S}$ operatorski sistem in $\mathfrak{B}$ $C^*$-algebra.
    Če je $\varphi: \mathcal{S} \to \mathfrak{B}$ pozitivna preslikava, potem je omejena in velja 
    $\| \varphi\| \leq 2 \cdot \varphi(1)$.
\end{trditev}

\begin{dokaz}
    Za sebiadjungiran $a \in \mathcal{S}$ velja
    $$-\| a\| \leq a \leq \| a\|$$
    in posledično $$-\varphi(\| a\|) \leq \varphi(a) \leq \varphi (\| a\|).$$
    Od tod sledi $| \varphi (a)| \leq \| a \| \cdot \varphi (1)$.
    Sedaj za poljuben $a \in \mathcal{S}$ dobimo 
    \begin{align*}
        \| \varphi(a)\| &= \| \varphi \left(\real a + i \cdot \imag a\right) \| \\
        &\leq \| \varphi (\real a)\| + \| \varphi (\imag a)\| \\
        &\leq (\| \real a\| + \| \imag a \|)\cdot \|\varphi (1)\| \\
        &\leq 2 \| a \| \cdot \| \varphi(1)\|. \qedhere
    \end{align*}
\end{dokaz}

Naslednji zgled (\cite[Appendix A.2, str. 221]{arveson1969}) pokaže, da je konstanta $2$ v oceni iz trditve \ref{trd:1.1} tudi dosežena.
\begin{zgled}[Arveson]\label{zgl:1.6}
    Naj bo $\mathcal{S} \subseteq C(\mathbb{T})$ operatorski sistem, ki ga razpenjajo funkcije $1, z, \overline{z}$.
    Definirajmo preslikavo
    $$\varphi: \mathcal{S} \to M_2 (\C), \quad \varphi (a + bz + c \overline{z}) = \begin{bmatrix}
        a & 2b\\
        2c & a
    \end{bmatrix}.$$
    Dokazujemo, da je $\varphi$ pozitivna preslikava.
    
    Trdimo, da velja 
    $$a + bz + c \overline{z} \geq 0 \quad \Leftrightarrow \quad \text{$b = \overline{c}$ in $a \geq 2|b|$}.$$
    Najprej implikacija v levo: če velja $b = \overline{c}$ in $a \geq 2|b|$, potem je 
    $$a + bz + c \overline{z} = a + 2\real{bz} \geq 0,$$
    saj $a \geq 2|b| \geq 2 \real{bz}$ za vse $z \in \mathbb{T}$.
    Nato še implikacija v desno: če je 
    $$\overline{a + bz + c \overline{z}} = a + bz + c \overline{z},$$
    sledi $a = \overline{a}$ in $c = \overline{b}$. Od tod dobimo
    $a + 2 \real (bz) \geq 0$ za vse $z \in \mathbb{T}$. Na krožnici obstaja tak $z \in \mathbb{T}$,
    da velja $bz = - |b|$, zato mora veljati tudi $a \geq 2 |b|$. 
    
    Če je $a + bz + \overline{bz} \geq 0$, potem se s 
    $\varphi$ preslika v matriko
    $$\begin{bmatrix}
        a & 2b\\
        2\overline{b} & a
    \end{bmatrix}.$$
    Slednja ima nenegativne diagonalne elemente in determinanto
    $$\det \begin{bmatrix}
        a & 2b\\
        2\overline{b} & a
    \end{bmatrix} = a^2 - 4 |b|^2 \geq 0,$$
    zato je pozitivna in $\varphi$ je pozitivna preslikava.
    
    Nazadnje opazimo, da velja
    $$\| \varphi \| \leq 2 \cdot \| \varphi(1) \| = 2 = \| \varphi(z)\| \leq \| \varphi\|,$$
    torej $\| \varphi \| = 2 \cdot \| \varphi(1)\|$.
\end{zgled}

\begin{opomba}
    V zgornjem zgledu sta dve težavi: kodomena $M_2 (\C)$ je nekomutativna (glej lemo \ref{lem:1.1})
    in operatorski sistem $\mathcal{S}$ `ne premore' dovolj pozitivnih elementov (glej trditev \ref{pos:russo-dye}).
    Zgled \ref{zgl:1.6} nam pove, da za posplošitev stanj na preslikave $\varphi: \mathfrak{A} \to \mathfrak{B}$
    ni dovolj predpostaviti zgolj pozitivnost, temveč potrebujemo povsem pozitivnost. 
\end{opomba}

Poglavje zaključimo s trditvijo, da je enotsko skrčitev na operatorskem prostoru $\mathcal{M}$
moč razširiti do pozitivne preslikave na operatorskem sistemu $\mathcal{M} + \mathcal{M}^*$.

\begin{trditev}
    Naj bo $\mathcal{S}$ operatorski sistem in $\mathfrak{B}$ $C^*$-algebra.
    Potem je vsaka enotska skrčitev $\varphi: \mathcal{S} \to \mathfrak{B}$ pozitivna.
\end{trditev}

\begin{dokaz}
    Po izreku GNS lahko (kot bomo v nadaljevanju pogosto storili) 
    brez škode za splošnost predpostavili, da je $\mathfrak{B} = \bh$ 
    za nek Hilbertov prostor $\mathcal{H}$.
    Fiksirajmo nek $x \in \mathcal{H}$, tako da velja $\| x\| = 1$.
    Oglejmo si funkcional
    $$\rho_x: \mathcal{S} \to \C,\quad \rho_x (a) = \langle \varphi(a) x, x \rangle.$$
    Po Cauchy-Schwarzevi neenakosti sledi
    $$|\rho_x (a)| \leq \| \varphi(a) x \| \cdot \| x\| \leq \| \varphi(a)\| \cdot \| x\|^2 = \|\varphi(a)\| \leq \| a\|,$$
    torej $\| \rho_x \| \leq 1$. Ker pa velja $\rho_x (1) = 1$, smo dokazali $\| \rho_x \| = 1 = \rho_x(1)$.
    Lema \ref{lem:1.1} nam zato pove, da je $\rho_x$ pozitiven funkcional.
    Sedaj vzemimo poljuben $a \in \mathcal{S}_+$. 
    Ker je
    $$\langle \varphi(a) x, x \rangle = \rho_x (a) \geq 0$$
    za vsak enotski vektor $x \in \mathcal{H}$, 
    mora $\varphi (a)$ biti pozitiven.
\end{dokaz}

\begin{trditev}
    Naj bo $\mathcal{M} \subseteq \mathfrak{A}$ operatorski prostor in $\mathfrak{B}$ $C^*$-algebra. 
    Potem lahko vsako enotsko skrčitev $\varphi: \mathcal{M} \to \mathfrak{B}$
    enolično razširimo do pozitivne preslikave na operatorskem sistemu $\mathcal{M} + \mathcal{M}^*$ s predpisom 
    $$\tilde{\varphi} : \mathcal{M} + \mathcal{M}^* \to \mathfrak{B},\quad \tilde{\varphi} (a + b^*) = \varphi(a) + \varphi (b)^*.$$
\end{trditev}

\begin{dokaz}
    Najprej dokažimo enoličnost. Če je $\psi: \mathcal{M} + \mathcal{M}^* \to \mathfrak{B}$ pozitivna razširitev
    $\varphi$, mora biti $*$-linearna. To pomeni, da za vsaka $a, b \in \mathcal{M}$ velja
    \begin{align*}
        \psi (a + b^*) &= \psi(a) + \psi(b^*) \\
        &= \psi(a) + \psi(b)^* \\
        &= \varphi(a) + \varphi(b)^* \\
        &= \tilde{\varphi} (a + b^*).
    \end{align*} 
    
    Nato dokažemo, da je $\tilde{\varphi}$ sploh dobro definirana preslikava.
    Vzemimo operatorski sistem 
    $$\mathcal{S} := \{a\ |\ \text{$a$ in $a^*$ sta elementa $\mathcal{M}$}\}.$$
    ($\mathcal{S}$ vsebuje enico, zato je neprazen).
    Ker je $\varphi: \mathcal{S} \to \mathfrak{B}$ enotska skrčitev, je pozitivna in v posebnem $*$-linearna.
    To pomeni, da za vsak $a \in \mathcal{S}$ velja
    $$\varphi (a^*) = \varphi(a)^*.$$
    Sedaj predpostavimo, da velja $a + b^* = c + d^* \in \mathcal{M} + \mathcal{M}^*$.
    Potem je $$a - c = (-b + d)^* \in \mathcal{S},$$
    torej 
    \begin{align*}
        \varphi(a) - \varphi(c) &= \varphi(a - c)\\
        &=\varphi(-b + d)^* \\
        &=-\varphi(b)^* + \varphi(d)^*.
    \end{align*}
    Od tod potem sledi
    \begin{align*}
        \tilde{\varphi}(a + b^*) &= \varphi(a) + \varphi(b)^*\\
        &= \varphi(c)  + \varphi(d)^*\\
        &= \tilde{\varphi}(c + d^*),
    \end{align*}
    torej je $\tilde{\varphi}$ dobro definirana.

    Nazadnje dokažimo, da je $\tilde{\varphi}$ pozitivna.
    Ponovno lahko predpostavimo, da je $\mathfrak{B} = \bh$ za nek Hilbertov prostor $\mathcal{H}$.
    Vzemimo nek $x \in \mathcal{H}$, za katerega velja $\| x\| = 1$. 
    Kot v prejšnji trditvi definiramo $$\rho_x : \mathcal{M} \to \C,\quad \rho_x (a) = \langle {\varphi} (a) x, x\rangle$$
    in 
    $$\tilde{\rho}_x : \mathcal{M} + \mathcal{M}^* \to \C,\quad \tilde{\rho}_x (a) = \langle \tilde{\varphi} (a) x, x\rangle.$$
    Dovolj je dokazati, da je $\tilde{\rho}_x$ pozitiven funkcional.
    Ker velja $ 1 = \rho_x (1) = \| \rho_x\|$, je $\rho_x$ po lemi \ref{lem:1.1} pozitiven funkcional. 
    Po Hahn-Banachovem izreku (\cite[Corollary III.6.5., str. 79]{conway1990}) ga lahko
    razširimo do funkcionala $\rho_x ': \mathcal{M} + \mathcal{M}^* \to \C$,
    za katerega prav tako velja $\| \rho_x '\| = 1 = \rho_x '(1)$.
    Slednji je zato pozitiven in posledično $*$-linearen.
    Od tod sledi, da za vsaka $a, b \in \mathcal{M}$ velja
    \begin{align*}
        \rho_x ' (a + b^*) = \rho_x (a) + \rho_x (b)^* = \tilde{\rho}_x (a + b^*).
    \end{align*}
    Torej je $\tilde{\rho}_x = \rho_x '$ pozitiven.
\end{dokaz}


\section{Povsem omejene preslikave}\label{sec:4}

Sedaj lahko postavimo zadostni pogoj za preslikavo $\varphi: \mathcal{M} \to \mathfrak{B}$,
da jo lahko razširimo do povsem pozitivne preslikave $\tilde{\varphi}: \mathcal{M} + \mathcal{M}^* \to \mathfrak{B}$.
Takšnim preslikavam pravimo povsem omejene preslikave. Ta razdelek je namenjen osnovnim lastnostim in zgledom povsem omejenih preslikav.

\begin{definicija}
    Naj bosta $\mathfrak{A}, \mathfrak{B}$ $C^*$-algebri in $\mathcal{M} \subseteq \mathfrak{A}$ operatorski prostor.
    Potem je preslikava $\varphi: \mathcal{M} \to \mathfrak{B}$:
    \begin{enumerate}
        \item \emph{povsem omejena}, če je $\| \varphi\|_{\po} := \sup_{n \in \N} \| \varphi^{(n)}\| < \infty$;
        \item \emph{povsem skrčitev}, če je $\| \varphi\|_{\po} \leq 1$;
        \item \emph{povsem izometrija}, če je $\varphi^{(n)}$ izometrija za vsak $n \in \N$. 
    \end{enumerate}
    Množico vseh povsem omejenih preslikav označujemo s $\PO (\mathcal{M}, \mathfrak{B})$.
\end{definicija}

\begin{opomba}
    Če je $n \geq m$, potem je $\| \varphi^{(n)} \| \leq \| \varphi^{(m)}\|$.
\end{opomba}
 
\begin{trditev}
    Množica $\PO (\mathcal{M}, \mathfrak{B})$ je Banachov prostor.
\end{trditev}

\begin{dokaz}
    Enostavno je videti, da je $\PO (\mathcal{M}, \mathfrak{B})$ vektorski prostor in $\| \cdot \|_{\po}$ res ustreza aksiomom za normo.
    Da dokažemo polnost, vzemimo poljubno Cauchyjevo zaporedje $\{\varphi_k\}_{k \in \N} \subseteq \PO (\mathcal{M}, \mathfrak{B})$.
    Ker je $\| \cdot \| \leq \| \cdot\|_{\po}$, je $\{\varphi_k\}$ Cauchyjevo tudi v navadni normi $\| \cdot\|$, torej $\varphi_k \to \psi$
    v navadni normi. Od tod sledi $\varphi_k ^{(n)} \to \psi^{(n)}$ za vsak $n \in \N$.
    Za poljuben $\varepsilon > 0$ obstaja tak $M \in \N$, da za vsaka $k, l \geq M$
    velja $\| \varphi_k - \varphi_l\|_{\po} < \varepsilon$, torej $\| \varphi_k ^{(n)} - \varphi_l ^{(n)}\|_{\po} < \varepsilon$ za vsak $n \in \N$.
    Če pošljemo $l \to \infty$, dobimo $\| \varphi_k ^{(n)} - \psi^{(n)} \| \leq \varepsilon$ za vsak $n \in \N$, od koder sledi 
    $\| \varphi_k - \psi\|_{\po} \leq \varepsilon$ za vse $k \geq M$.
    To pomeni, da $\varphi_k \to \psi$ tudi v normi $\| \cdot \|_{\po}$.    
\end{dokaz}

\begin{trditev}
    Naj bo $\mathcal{S}$ operatorski sistem in $\mathfrak{B}$ $C^*$-algebra.
    Potem je vsaka enotska povsem skrčitev $\varphi: \mathcal{S} \to \mathfrak{B}$ povsem pozitivna.
\end{trditev}

\begin{dokaz}
    Za vsak $n \in \N$ je $\varphi^{(n)}: M_n (\mathcal{S}) \to M_n (\mathfrak{B})$ enotska skrčitev,
    zato je pozitivna.
\end{dokaz}

\begin{posledica}\label{pos:1.2}
    Vsaka enotska povsem skrčitev $\varphi: \mathcal{M} \to \mathfrak{B}$ premore enolično določeno povsem pozitivno razširitev 
    $\tilde{\varphi}: \mathcal{M} + \mathcal{M}^* \to \mathfrak{B}$.
\end{posledica}

\begin{dokaz}
    Za poljuben $n \in \N$ je $\varphi^{(n)}: M_n (\mathcal{M}) \to M_n (\mathfrak{B})$ enotska skrčitev, 
    torej je $\widetilde{(\varphi^{(n)})}: M_n (\mathcal{M}) + M_n (\mathcal{M})^* \to M_n (\mathfrak{B})$ pozitivna.
    Ker je $M_n (\mathcal{M}) + M_n (\mathcal{M})^* = M_n (\mathcal{M} + \mathcal{M}^*)$ in preslikava
    $\widetilde{(\varphi^{(n)})}$ sovpada s $(\tilde{\varphi})^{(n)}: M_n (\mathcal{M} + \mathcal{M}^*) \to M_n (\mathfrak{B})$,
    je slednja preslikava prav tako pozitivna.
    \begin{comment}
    Vzemimo matriki $[a_{ij}]_{i, j} \in M_n (\mathcal{M}), [b_{ij} ^*]_{i, j} \in M_n (\mathcal{M})$. Potem je 
    \begin{align*}
        \widetilde{(\varphi^{(n)})} ([a_{ij}]_{i, j} + [b_{ij} ^*]_{i, j}) &= \varphi^{(n)} ([a_{ij}]_{i, j}) + \varphi^{(n)} ([b_{ji}]_{i, j})^*\\
        &= [\varphi(a_{ij})]_{i, j} + [\varphi(b_{ji})]_{i, j}^*\\
        &= [\varphi(a_{ij})]_{i, j} + [\varphi(b_{ij})^*]_{i, j}\\
        &= [\varphi(a_{ij})]_{i, j} + [\varphi(b_{ij}^*)]_{i, j}\\
        &= [\varphi(a_{ij} + b_{ij}^*)]_{i, j}\\
        &= [\tilde{\varphi} (a_{ij} + b_{ij}^*)]_{i,j}\\
        &= \tilde{\varphi}^{(n)} ([ a_{ij} + b_{ij}^*]_{i,j})\\
        &= \tilde{\varphi}^{(n)} ([ a_{ij}]_{i, j} + [b_{ij}^*]_{i,j}),
    \end{align*}
    torej je $\tilde{\varphi}^{(n)}$ pozitivna za vsak $n \in \N$.
    \end{comment}
\end{dokaz}

\begin{opomba}
    Z oznako $\os$ bomo označevali kategorijo operatorskih prostorov, v kateri so morfizmi povsem omejene preslikave, 
    v $\os_1$ pa enotske povsem omejene preslikave.
\end{opomba}

Oglejmo si nekatere zglede povsem omejenih preslikav.

\begin{zgled}
    Dokažimo, da je vsak omejen funkcional $\varphi: \mathcal{M} \to \C$ povsem omejen.
    Za poljubna enotska vektorja $\alpha, \beta \in \C^n$ in $A \in M_m (\mathcal{M})$ imamo
    \begin{align*}
        \left| \beta^* \varphi^{(n)} (A) \alpha \right| &= \left| \varphi(\beta^* A \alpha)\right|\\
        &\leq \| \varphi\| \cdot \| \beta^* A \alpha\|\\
        &\leq \| \varphi\| \cdot \| A\| \cdot \| \alpha \| \cdot \| \beta\|\\
        &= \| \varphi\| \cdot \| A\|.
    \end{align*}
    Če vzamemo supremum obeh strani po $\alpha, \beta$, dobimo 
    $$\| \varphi^{(n)} (A)\| \leq \| \varphi\| \cdot \| A\|$$
    in posledično $\| \varphi^{(n)} \| = \| \varphi\|= \| \varphi \|_{\po}$.
\end{zgled}

\begin{comment}
\begin{zgled}
    Iz zgleda \ref{zgl:1.3} že vemo, da je vsak pozitivni funkcional $\varphi: \mathcal{S} \to \C$ povsem pozitiven.
    Dokazujemo, da je povsem omejena, in sicer da velja $\| \varphi^{(n)}\| = \| \varphi\|$ za vsak $n \in \N$.
    Brez škode za splošnost predpostavimo, da je $\| \varphi \| = \varphi(1) = 1$.
    Dovolj je dokazati Cauchy-Schwarzevo neenakost
    \begin{equation}
        \varphi^{(n)} (A) ^* \varphi^{(n)}(A) \leq \varphi^{(n)} (A^* A) \label{eq:1.3}
    \end{equation}
    za poljubno matriko $A \in M_n (\mathcal{S})$. Od tod po pozitivnosti $\varphi^{(n)}$ sledi 
    $$\|\varphi^{(n)} (A)\|^2 = \| \varphi^{(n)} (A)^* \varphi^{(n)} (A) \| \leq \| \varphi^{(n)} (A^* A)\| \leq \| \varphi^{(n)} (1)\| \cdot \| A^* A\| = \| A\|^2,$$
    torej $\| \varphi^{(n)}\| = \| \varphi\| = 1$ za vsak $n \in \N$. Preostane le še, da dokažemo neenakost \eqref{eq:1.3}.
    Vzemimo poljuben vektor $\alpha = (\alpha_1, \dots, \alpha_n) \in \C^n$. Potem je po uporabi neenakosti \eqref{eq:1.2}
    \begin{align*}
        \langle \varphi^{(n)} (A)^* \varphi^{(n)} (A) \alpha, \alpha \rangle &= \langle \varphi^{(n)} (A) \alpha, \varphi^{(n)} (A) \alpha \rangle = \sum_{i = 1} ^n \left| \sum_{k = 1} ^n \varphi(a_{ik}) \alpha_k \right|^2\\
        &= \sum_{i = 1} ^n \left| \varphi\left(\sum_{k = 1} ^n a_{ik} \alpha_k \right) \right|^2 \leq \sum_{i = 1} ^n \varphi\left(\left(\sum_{k = 1} ^n a_{ik} \alpha_k \right)^* \left(\sum_{j = 1} ^n a_{ij} \alpha_j \right)\right) \\
        &= \sum_{i = 1} ^n \varphi\left(\sum_{k = 1} ^n a_{ik}^* \overline{\alpha_k} \sum_{j = 1} ^n a_{ij} \alpha_j \right) = \sum_{i = 1} ^n \sum_{k = 1} ^n \sum_{j = 1} ^n \alpha_j \varphi\left( a_{ik}^* a_{ij} \right)  \overline{\alpha_k} \\
        &= \sum_{k = 1} ^n \sum_{j = 1} ^n \sum_{i = 1} ^n  \alpha_j \varphi\left( a_{ik}^* a_{ij} \right)  \overline{\alpha_k} = \sum_{k = 1} ^n \sum_{j = 1} ^n  \alpha_j \varphi\left(\sum_{i = 1} ^n  a_{ik}^* a_{ij} \right)  \overline{\alpha_k} \\
        &= \langle \varphi^{(n)} (A^* A) \alpha, \alpha \rangle, 
    \end{align*}
    torej $\varphi^{(n)} (A)^* \varphi^{(n)} (A) \leq \varphi^{(n)} (A^* A)$.
\end{zgled}
\end{comment}

\begin{zgled}
    Vemo že, da je vsak enotski $*$-homomorfizem $\varphi: \mathfrak{A} \to \mathfrak{B}$ povsem pozitiven.
    Velja pa tudi, da je $\varphi$ skrčitev. Ker so preslikave $\varphi^{(n)}$ prav tako enotski $*$-homomorfizmi,
    so tudi skrčitve in je $\varphi$ povsem skrčitev. Če je $\varphi$ enotski $*$-monomorfizem, je izometrija in enako velja za vse njegove ampliacije $\varphi^{(n)}$.
    Torej je vsak enotski $*$-monomorfizem med $C^*$-algebrama povsem izometrija.
\end{zgled}

Opazimo, da so vse povsem pozitivne preslikave, ki jih do sedaj poznamo (pozitivni funkcionali, $*$-homomorfizmi itd.), tudi povsem omejene.
V preostanku tega poglavja dokažemo, da to velja tudi v splošnem. Argumentom, ki sledijo, pravimo tudi \emph{dilatacijski triki}.

\begin{lema}\label{lem:random1}
    Naj bo $\mathcal{H}$ Hilbertov prostor in $P, Q, A \in \bh$.
    Potem je $$\begin{bmatrix}
            P & A\\
             A^* & Q
        \end{bmatrix} \geq 0$$
        natanko tedaj, ko sta $P, Q$ pozitivna operatorja in 
        $|\langle Ax, y\rangle|^2 \leq \langle Py, y\rangle \cdot \langle Qx, x \rangle$
        za vse $x, y \in \mathcal{H}$.
        V posebnem velja
        $$\begin{bmatrix}
            1 & A\\
             A^* & Q
        \end{bmatrix} \geq 0$$
        natanko tedaj, ko $A^* A \leq Q$.
\end{lema}

\begin{dokaz}
    Predpostavimo, da velja $|\langle Ax, y\rangle|^2 \leq \langle Py, y\rangle \cdot \langle Qx, x \rangle$ za vse $x, y \in \mathcal{H}$.
        Potem sledi
        \begin{align*}
            \left\langle \begin{bmatrix}
                P & A\\
                 A^* & Q
            \end{bmatrix} \begin{bmatrix}
                y\\ x
            \end{bmatrix}, \begin{bmatrix}
                y \\ x
            \end{bmatrix} \right\rangle &= \langle Py, y \rangle + 2\real \langle Ax, y \rangle + \langle Qx, x\rangle\\
            &\geq \langle Py, y \rangle - 2 |\langle Ax, y \rangle| + \langle Qx, x\rangle\\
            &\geq \langle Py, y \rangle - 2 \langle P y, y \rangle^{\frac{1}{2}} \langle Q x, x \rangle^{\frac{1}{2}}  + \langle Qx, x\rangle\\
            &= (\langle Py, y \rangle^{\frac{1}{2}} - \langle Qx, x \rangle^{\frac{1}{2}})^2 \geq 0.
        \end{align*}
        Predpostavimo še obratno, torej da obstajata $x, y$, tako da
        $|\langle Ax, y\rangle|^2 > \langle Py, y\rangle \cdot \langle Qx, x \rangle$.
        Brez škode za splošnost smemo predpostaviti, da je $\langle P y, y \rangle = \langle Qx, x \rangle = 1$, 
        saj bi v nasprotnem primeru lahko pomnožili $x, y$ z ustreznima faktorjema.
        Prav tako lahko predpostavimo $\langle Ax, y \rangle \leq 0$,
        saj lahko v nasprotnem primeru preprosto zavrtimo $x$ za ustrezen kot $e^{i \theta}$. 
        Od tod sledi $$\langle Ax, y \rangle = - |\langle Ax, y \rangle | < -1$$
        in 
        \begin{align*}
            \left\langle \begin{bmatrix}
                P & A\\
                 A^* & Q
            \end{bmatrix} \begin{bmatrix}
                y\\ x
            \end{bmatrix}, \begin{bmatrix}
                y \\ x
            \end{bmatrix} \right\rangle &= \langle Py, y \rangle + 2\real \langle Ax, y \rangle + \langle Qx, x\rangle\\
            &< 1 - 2 + 1 = 0.
        \end{align*}
        Za dokaz druge trditve je dovolj pokazati,
        da velja $|\langle Ax, y\rangle|^2 \leq \langle y, y\rangle \cdot \langle Qx, x \rangle$ za vse $x, y$.
        Če vstavimo $y = Ax$, dobimo
        $$\langle Ax, Ax \rangle \leq \langle Qx, x \rangle$$ za vse $x$.
        Ker je leva stran enačbe enaka $\langle A^* A x, x \rangle$, dobimo $A^* A \leq Q$.
        Obratno, če velja $A^* A \leq Q$, potem po Cauchy-Schwarzove neenakosti sledi
        \begin{align*}
            |\langle Ax, y \rangle|^2 &\leq \| Ax \|^2 \cdot \| y\|^2 \\
            &\leq \langle A^* A x, x \rangle \cdot \langle y, y \rangle \\
            &\leq \langle Qx, x \rangle \cdot \langle y, y \rangle. \qedhere
        \end{align*}
\end{dokaz}

\begin{posledica}
    Če je $$\begin{bmatrix}
            P & A\\
             A^* & Q
        \end{bmatrix} \geq 0,$$
    potem velja $A^* A \leq \| P\| Q$ in $A^* A \leq \| Q\| P$.
\end{posledica}

\begin{dokaz}
    Ker je $0 \leq P \leq \| P\|$, imamo
        $$\begin{bmatrix}
            \| P\| & A\\
             A^* & Q
        \end{bmatrix} \geq \begin{bmatrix}
            P & A\\
             A^* & Q
        \end{bmatrix} \geq 0.$$
        Od tod sledi 
        $$\begin{bmatrix}
            1 & \frac{A}{\| P\|}\\
            \frac{A^*}{\| P\|} &\frac{Q}{\| P\|}
        \end{bmatrix} \geq 0$$ in 
        $$\frac{A^*}{\| P\|} \frac{A}{\| P\|} \leq \frac{Q}{\| P\|}$$
        po prejšnji lemi. Od tod sledi $A^* A \leq \| P\| Q$.
        Po enakem argumentu dokažemo tudi drugo neenakost.
\end{dokaz}

\begin{posledica}\label{pos:1.3}
    Naj bo $\mathcal{S}$ operatorski sistem in $\varphi: \mathcal{S} \to \mathfrak{B}$ $2$-pozitivna preslikava.
    Potem je $\| \varphi \| = \| \varphi(1)\|$.
\end{posledica}

\begin{dokaz}
    Ker je $\varphi$ $1$-pozitivna, je tudi $*$-linearna.
    Vzemimo $a \in \mathcal{S}$, tako da je $\| a\| \leq 1$.
    Potem je $$\begin{bmatrix}
        1 & a\\
        a^* & 1
    \end{bmatrix} \geq 0,$$
    torej 
    $$\begin{bmatrix}
        \varphi(1) & \varphi(a)\\
        \varphi(a)^* & \varphi(1)
    \end{bmatrix} \geq 0.$$
    Po prejšnji posledici dobimo $$\varphi (a) ^* \varphi(a) \leq \| \varphi(1)\| \varphi(1)$$
    in zato $\| \varphi(a)\|^2 \leq \| \varphi(1)\|^2$.
    S tem smo dokazali $\| \varphi \| \leq \| \varphi(1)\|$.
\end{dokaz}

\begin{posledica}[Cauchy-Schwarz za $2$-pozitivne preslikave]
    Naj bosta $\mathfrak{A}, \mathfrak{B}$ $C^*$-algebri in $\varphi: \mathfrak{A} \to \mathfrak{B}$ 
    $2$-pozitivna preslikava. Potem velja $$\varphi (a^* b) \varphi (b^* a) \leq  \| \varphi (a^* a ) \| \cdot \varphi (b^* b)$$
    in v posebnem
    $$\varphi (a)^* \varphi (a) \leq \| \varphi(1)\| \varphi(a^* a).$$
\end{posledica}

\begin{dokaz}
    Ker je
    $$\begin{bmatrix}
        a & b\\
        0 & 0
    \end{bmatrix}^* \begin{bmatrix}
        a & b\\
        0 & 0
    \end{bmatrix} = \begin{bmatrix}
        a^* a & a^* b\\
        b^* a & b^* b
    \end{bmatrix} \geq 0,$$
    dobimo
    $$\begin{bmatrix}
        \varphi(a^* a) & \varphi(a^* b)\\
        \varphi(a^* b)^* & \varphi(b^* b)
    \end{bmatrix} \geq 0$$
    in posledično $\varphi (a^* b) \varphi(b a^*) \leq \| \varphi(a^* a) \| \varphi(b^* b)$.
\end{dokaz}

Sedaj lahko zapišemo glavno trditev tega poglavja.

\begin{trditev}
    Naj bo $\mathcal{S} \subseteq \mathfrak{A}$ operatorski sistem in $\mathfrak{B}$ $C^*$-algebra.
    Če je $\varphi: \mathcal{S} \to \mathfrak{B}$ povsem pozitivna, potem je tudi povsem omejena in velja
    $\| \varphi \| = \| \varphi \|_{\po} = \| \varphi (1)\|$.
\end{trditev}

\begin{dokaz}
    Ker je $\| \varphi (1) \| \leq \| \varphi\| \leq \| \varphi\|_{\po}$, zadošča pokazati
    $\| \varphi\|_{\po} \leq \| \varphi(1)\|$. 
    Vzemimo $n \in \N$ in $a \in M_n (\mathcal{S})$,
    tako da velja $\| a\| = 1$. Ker je $\varphi$ povsem pozitivna preslikava, je
    $$ \varphi^{(2n)} : M_2 (M_n (\mathcal{S})) \to M_2 (M_n (\mathcal{S}))$$
    pozitivna. 
    Ker pa velja $ \varphi^{(2n)} = (\varphi^{(n)})^{(2)}$, je $\varphi^{(n)}$ $2$-pozitivna preslikava.
    Po posledici \ref{pos:1.3} nato sledi $\| \varphi^{(n)}\| = \| \varphi^{(n)} (I_n)\| = \| \varphi (1)\|$.
\end{dokaz}

\begin{opomba}
    V prejšnjem dokazu smo identificirali $M_{2n} \equiv M_2 (M_n(\mathcal{S}))$.
    To smemo narediti, saj imamo za poljubna $m, n \in \N$ $*$-izomorfizem
    $$\Phi_{m, n} : M_m (M_n (\mathfrak{A})) \to M_{mn} (\mathfrak{A}),\quad \begin{bmatrix}
        [A_{11}] & \cdots & [A_{1m}]\\
        \vdots & & \vdots\\
        [A_{m1}] & \cdots & [A_{mm}]
    \end{bmatrix} \mapsto \begin{bmatrix}
        A_{11} & \cdots & A_{1m}\\
        \vdots & & \vdots\\
        A_{m1} & \cdots & A_{mm}
    \end{bmatrix},$$
    ki ohranja pozitivne elemente. 
\end{opomba}

\begin{posledica}
    Če je $\mathcal{M} \subseteq \mathfrak{A}$ operatorski prostor in $\varphi: \mathcal{M} \to \mathfrak{B}$
    enotska povsem izometrija, potem je tudi $\tilde{\varphi}: \mathcal{M} + \mathcal{M}^* \to \mathfrak{B}$
    povsem izometrija.
\end{posledica}

\begin{dokaz}
    Ker je $\varphi: \mathcal{M} \to \mathfrak{B}$ enotska skrčitev, je tudi 
    $$\tilde{\varphi}: \mathcal{M} + \mathcal{M}^* \to \varphi(\mathcal{M}) + \varphi(\mathcal{M})^*$$
    povsem skrčitev. Po drugi strani pa je tudi 
    $\varphi^{-1}: \varphi(\mathcal{M}) \to \mathfrak{A}$ enotska skrčitev, zato je
    $$\widetilde{\varphi^{-1}}: \varphi(\mathcal{M}) + \varphi(\mathcal{M})^* \to \mathcal{M} + \mathcal{M}^*$$
    povsem skrčitev. Ker sta preslikavi $\tilde{\varphi}$ in $\widetilde{\varphi^{-1}}$
    druga drugi inverzni in obe povsem skrčitvi, sta obe povsem izometrični.
\end{dokaz}

\begin{posledica}
    Množica $\PP (\mathcal{S}, \mathfrak{B})$ je konveksen stožec v $\PO (\mathcal{S}, \mathfrak{B})$.
\end{posledica}

Naslednja trditev je posplošitev leme \ref{lem:1.1} za povsem pozitivne preslikave.

\begin{posledica}
    Naj bo $\varphi: \mathcal{S} \to \mathfrak{B}$ enotska preslikava.
    Potem je $\varphi$ povsem pozitivna natanko tedaj, ko je $\| \varphi \|_{\po} = \| \varphi\| = 1$.
\end{posledica}

\begin{dokaz}
    Sledi direktno iz posledice \ref{pos:1.2} in prejšnje trditve.
\end{dokaz}

\begin{comment}

\begin{lema}
    Če je $\mathcal{M}$ operatorski prostor in $\varphi: \mathcal{M} \to M_n (\C)$ omejena preslikava,
    potem je povsem omejena in velja $\| \varphi\|_{\po} \leq n \| \varphi\|$.
\end{lema}

\begin{dokaz}
    Najprej dokažimo trditev za primer $n = 1$. Tedaj imamo za poljubna vektorja $\alpha, \beta \in \C^n$ in $A \in M_m (\mathcal{M})$
    \begin{align*}
        \left| \langle \varphi^{(m)} (A) \alpha, \beta \rangle \right| &= \left| \varphi(\beta^\top A \alpha)\right|\\
        &\leq \| \varphi\| \cdot \| \beta^\top A \alpha\|\\
        &\leq \| \varphi\| \cdot \| A\| \cdot \| \alpha \| \cdot \| \beta\|.
    \end{align*}
    Če vzamemo supremum obeh strani neenačbe po enotskih vektorjih $\| \alpha\| = \| \beta \| = 1$, dobimo 
    $$\| \varphi^{(m)} (A)\| \leq \| \varphi\| \cdot \| A\|$$
    in posledično $\| \varphi \|_{\po} = \| \varphi\|$.
    Sedaj dokažimo trditev za splošen $n \in \N$.
    Definirajmo preslikave $$\varphi_{ij}: \mathcal{M} \to \C,\quad \varphi_{ij} (a) = \langle \varphi(a) e_j, e_i \rangle.$$
    Ker je $\varphi (a) = \sum_{i, j} ^n \varphi_{ij} (a) E_{ij}$, lahko zapišemo njene ampliacije kot  
    $$\varphi^{(m)} (A) = \sum_{i, j} ^n \varphi_{ij}^{(m)} (A) \otimes E_{ij} \in M_m (M_n (\C)).$$
    Opazimo tudi, da je $$\Phi: M_n (\C) \otimes M_m (\C) \to M_m (\C) \otimes M_n (\C),\quad A \otimes B \mapsto B \otimes A$$
    $*$-izomorfizem $C^*$-algeber, zato ohranja normo.
    Od tod sledi 
    \begin{align*}
        \| \varphi^{(m)} (A)\| &= \left\| \sum_{i, j} ^n \varphi_{ij}^{(m)} (A) \otimes E_{ij} \right\|\\
        &=  \left\|\sum_{i, j} ^n E_{ij} \otimes \varphi_{ij}^{(m)} (A) \right\|\\
        &= \left\| \begin{bmatrix}
            \varphi_{11} ^{(m)} (A) & \cdots & \varphi_{1n} ^{(m)} (A)\\
            \vdots & & \vdots\\
            \varphi_{n1} ^{(m)} (A) & \cdots & \varphi_{nn} ^{(m)} (A)\\
        \end{bmatrix} \right\|\\
        &= \left\| \begin{bmatrix}
            \left\| \varphi_{11} ^{(m)} (A) \right\| & \cdots & \left\| \varphi_{1n} ^{(m)} (A) \right\|\\
            \vdots & & \vdots\\
            \left\| \varphi_{n1} ^{(m)} (A) \right\| & \cdots & \left\| \varphi_{nn} ^{(m)} (A) \right\|\\
        \end{bmatrix} \right\|\\
        &\leq \left\| \begin{bmatrix}
            \left\| \varphi_{11} ^{(m)} (A) \right\| & \cdots & \left\| \varphi_{1n} ^{(m)} (A) \right\|\\
            \vdots & & \vdots\\
            \left\| \varphi_{n1} ^{(m)} (A) \right\| & \cdots & \left\| \varphi_{nn} ^{(m)} (A) \right\|\\
        \end{bmatrix} \right\|_{F}\\
        &\leq n \max_{ij} \| \varphi_{ij} ^{(m)} (A)\|\\
        &\leq n \cdot \| \varphi \| \cdot \| A\|,
    \end{align*}
    pri čemer $\| \cdot \|_F$ označuje Frobeniusovo matrično normo.
\end{dokaz}



\begin{zgled}
    Preslikava $$\varphi: M_2 (\C) \to M_2 (\C),\quad A \mapsto A^\top$$
    iz zgleda \ref{zgl:1.4} je po prejšnji lemi povsem omejena, toda ni povsem pozitivna.
\end{zgled}
\end{comment}

\section{Povsem pozitivne preslikave na $C(X)$}\label{sec:5}


V prejšnjih poglavjih smo že srečali primere povsem pozitivnih preslikav 
$\varphi: \mathfrak{A} \to \mathfrak{B}$, kot na primer pozitivne funkcionale, $*$-homomorfizme
in enotske povsem skrčitve. V tem poglavju se omejimo na primer, ko je vsaj ena izmer $C^*$-algeber $\mathfrak{A}, \mathfrak{B}$
komutativna. Po Gelfandovi transformaciji (\cite[Theorem VIII.2.1., str. 236]{conway1990}) je vsaka komutativna $C^*$-algebra
$*$-izomorfna $C(X)$, kjer je $X$ nek kompakten Hausdorffov topološki prostor (če ni izrecno napisano drugače,
bomo v tem magistrskem delu za $C(X)$ vedno implicitno predpostavili, da je $X$ kompakten ter Hausdorfov).
Stinespring (\cite{stinespring1955}) je dokazal, da za povsem pozitivne preslikave
$\varphi: C(X) \to \mathfrak{B}$ in $\varphi: \mathfrak{A} \to C(X)$ obstaja preprosta karakterizacija.

\begin{trditev}\label{trd:1.2}
    Preslikava $\varphi: C(X) \to \mathfrak{B}$ je povsem pozitivna natanko tedaj, ko je pozitivna. 
\end{trditev}

\begin{dokaz}
    Brez škode za splošnost predpostavimo, da je $\mathfrak{B} = \bh$ za nek Hilbertov prostor $\mathcal{H}$.
            Fiksirajmo neka vektorja $u, v\in \mathcal{H}$. Po Rieszovem reprezentacijskem izreku \cite[Theorem 6.18, str. 129]{rudin1987} obstaja (enolična) 
            regularna kompleksna Borelova mera $\mu_{u, v}$,
            tako da velja
            $$\langle \varphi(a) v, u \rangle = \int_X a \, d\mu_{u, v},\quad \forall a \in \mathfrak{A}.$$
            Vzemimo poljubnih $n$ vektorjev $u_1, \dots, u_n \in \mathcal{H}$ in definiramo pozitivno mero
            $\mu := \sum_{i, j} ^n |\mu_{u_i, u_j}|$. Označimo s
            $h_{u_i, u_j} = \frac{d \mu_{u_i, u_j}}{d \mu}$ Radon--Nikodymove odvode. 
            Dokazujemo, da je matrika $$H(x) = [h_{u_i, u_j} (x)]_{i, j}$$
            pozitivna $\mu$-skoraj povsod.    
            Vzemimo poljubne skalarje $\lambda_1, \dots, \lambda_n \in \C$ in definiramo $u := \sum_{i = 1} ^n \lambda_i u_i$.
            Tedaj je
            \begin{align*}
                \langle \varphi (a) u, u \rangle &= \left\langle \varphi (a) \left( \sum_{j = 1} ^n \lambda_j u_j \right), \sum_{i = 1} ^n \lambda_i u_i \right\rangle\\
                &= \sum_{i, j} ^n \overline{\lambda_i} \lambda_j \langle \varphi (a) u_j, u_i \rangle\\
                &= \sum_{i, j} ^n \overline{\lambda_i} \lambda_j \int_X a\, d\mu_{u_i, u_j}.
            \end{align*}
            Ker je $\varphi$ pozitivna preslikava, je
            $$\mu_{u, u} = \sum_{i, j} ^n \overline{\lambda_i} \lambda_j \mu_{u_i, u_j}$$
            pozitivna mera.
            Opazimo, da velja $\mu_{u, u} \ll \mu$, torej je njen Radon--Nikodymov odvod
            $$\frac{ d\mu_{u, u}}{d \mu} = \sum_{i, j} ^n \overline{\lambda_i} \lambda_j h_{u_i, u_j} \geq 0$$
            $\mu$-skoraj povsod.
            Dokazali smo, da je za vsak $\lambda \in \C^n$ funkcija $\langle H(x) \lambda, \lambda \rangle$ nenegativna $\mu$-skoraj povsod.
            Vzemimo nek gost števen podprostor $D \subseteq \C^n$, recimo $D = (\Q + i \Q)^n$.
            Za vsak $\lambda \in D$ označimo $$N_{\lambda} := \{x \in X\ |\ \langle H(x) \lambda, \lambda \rangle < 0\}$$
            ter $N := \bigcup_{\lambda \in D} N_\lambda$.
            Ker je $\mu(N_\lambda) = 0$ za vsak $\lambda \in D$ in je $D$ števna množica, je tudi $\mu(N) = 0$.
            Vzemimo poljuben $x \notin N$.
            Potem je zvezna funkcija $\lambda \mapsto \langle H(x) \lambda, \lambda \rangle$ nenegativna na $D \subseteq \C^n$, torej je nenegativna na celotnem
            prostoru $\C^n$. To pomeni, da je $H(x)$ nenegativna matrika povsod razen na $\mu$-ničelni množici $N$.
            Sedaj vzemimo poljubne $n \in \N$, $a_i, \dots, a_n \in \mathfrak{A}$ in $u_1, \dots, u_n \in \mathcal{H}$.
            Potem velja
            \begin{align*}
                \left\langle \varphi^{(n)} \left(\begin{bmatrix}
                    a_1^* a_1  & \cdots & a_1 ^* a_n\\
                    \vdots & & \vdots \\
                    a_n ^* a_1  & \cdots & a_n ^* a_n\\
                \end{bmatrix}\right) \begin{bmatrix}
                    u_1 \\ \vdots \\ u_n
                \end{bmatrix}, 
                \begin{bmatrix}
                    u_1 \\ \vdots \\ u_n
                \end{bmatrix}
                \right\rangle &= \sum_{i, j} ^n \langle \varphi (a_i ^* a_j) u_j, u_i \rangle \\
                &= \sum_{i, j} ^n \int_X \overline{a_i} (x) a_j (x)\, d\mu_{u_i, u_j}\\
                &= \sum_{i, j} ^n \int_X \overline{a_i} (x) a_j (x) h_{u_i, u_j}(x)\, d\mu\\
                &= \int_X \sum_{i, j} ^n \overline{a_i} (x) a_j (x) h_{u_i, u_j}(x)\, d\mu \geq 0\\ 
                &= \int_X \left\langle H(x) \begin{bmatrix}
                    a_1 (x) \\ \vdots \\ a_n (x)
                \end{bmatrix}, \begin{bmatrix}
                    a_1 (x) \\ \vdots \\ a_n (x)
                \end{bmatrix} \right\rangle\, d\mu\\ 
                &\geq 0.  \qedhere
            \end{align*}
\end{dokaz}

\begin{opomba}
    Če je $\mathcal{S} \subseteq C(X)$ operatorski sistem, potem pozitivna preslikava $\varphi: \mathcal{S} \to \mathfrak{B}$
    še ni nujno povsem pozitivna.
    Protiprimer smo že srečali v zgledu \ref{zgl:1.6}, kjer smo definirali 
    pozitivno preslikavo $\varphi: \mathcal{S} \to M_2 (\C)$ na operatorskem sistemu $\mathcal{S} = \lin\{1, z, \overline{z}\} \subseteq C(\T)$,
    tako da $\varphi$ ni povsem pozitivna.
    Težava je v tem, da razcep iz leme \ref{lem:1.2} ne velja za poljuben operatorski sistem, temveč zgolj za celotno $C^*$-algebro $\mathfrak{A}$.
\end{opomba}

Sedaj dokažimo podobno trditev še za primer, ko je kodomena $C(X)$.

\begin{trditev}
    Naj bo $\mathcal{S} \subseteq \mathfrak{A}$ poljuben operatorski sistem.
    Potem je preslikava $\varphi: \mathcal{S} \to C(X)$ povsem pozitivna natanko tedaj, ko je pozitivna.
\end{trditev}


\begin{dokaz}
    \begin{enumerate}
        Vzemimo poljubno pozitivno matriko $A = [a_{ij}] \in M_n(\mathcal{S})$ in $b_1, \dots, b_n \in C(X)$. Dokazati želimo, da velja
        $\sum_{i, j} ^n b_j^* \varphi(a_{ij}) b_i \geq 0$. Za poljuben $x \in X$ imamo
        \begin{align*}
            \left(\sum_{i, j} ^n b_j^* \varphi(a_{ij}) b_i \right) (x) &= \sum_{i, j} ^n \overline{b_j} (x) \cdot \varphi(a_{ij}) (x) \cdot b_i (x)\\
            &= \sum_{i, j} ^n  \varphi(\overline{b_j (x)} \cdot a_{ij} \cdot b_i (x)) (x) \\
            &=  \varphi \left(\sum_{i, j} ^n \overline{b_j (x)} \cdot a_{ij} \cdot b_i (x) \right) (x) \\
            &= \varphi \left( \begin{bmatrix}
                b_1 (x) \\ \vdots \\ b_n (x)
            \end{bmatrix}^* A \begin{bmatrix}
                b_1 (x) \\ \vdots \\ b_n (x)
            \end{bmatrix} \right) (x) \geq 0. \qedhere
        \end{align*}
    \end{enumerate}
\end{dokaz}

\begin{posledica}
    Če je $\varphi: C(X) \to \mathfrak{B}$ enotska skrčitev, potem je povsem skrčitev.
    Enako velja tudi za enotsko skrčitev $\varphi: \mathcal{M} \to C(X)$,
    kjer je $\mathcal{M} \subseteq \mathfrak{A}$ poljuben operatorski prostor.
\end{posledica}

\begin{dokaz}
    Ker je $\varphi: C(X) \to \mathfrak{B}$ enotska skrčitev, je pozitivna.
    Po trditvi \ref{trd:1.2} je tudi povsem pozitivna in zato povsem skrčitev. 
    Enak argument velja tudi za enotsko skrčitev $\varphi: \mathcal{M} \to C(X)$.
\end{dokaz}

Naslednjo posledico bomo večkrat uporabili kasneje.

\begin{posledica}\label{pos:aux}
    Naj bo $\mathcal{S} \subseteq C(X)$ operatorski sistem.
    Če je $\varphi: \mathcal{S} \to C(Y)$ enotska izometrija, potem je povsem izometrija.
\end{posledica}

\begin{dokaz}
    Ker je $\varphi: \mathcal{S} \to C(X)$ enotska skrčitev, je po prejšnji posledici povsem skrčitev.
    Po drugi strani pa je njen inverz $\varphi^{-1}: \varphi(\mathcal{S}) \to C(X)$ prav tako enotska skrčitev,
    zato je povsem skrčitev. Ker sta tako $\varphi$ in $\varphi^{-1}$ povsem skrčitvi, je $\varphi$ povsem izometrija.
\end{dokaz}

\section{Dilatacije}\label{sec:6}

Dilatacije so orodje, s katerimi operator na Hilbertovem prostoru razširimo do `pohlevnejšega' operatorja na večjem Hilbertovem prostoru.
To nam omogoča, da lahko nekatere trditve za splošne operatorje reduciramo na bolj poseben razred operatorjev. 
V tem razdelku dokažemo Stinespringov izrek (\cite{stinespring1955}), ki dilatacije poveže s povsem pozitivnimi preslikavami.


\begin{definicija}
    \emph{Dilatacija} operatorja $T \in \bh$ je operator $S \in \bk$, pri čemer 
    je $\mathcal{H}$ podprostor Hilbertovega prostora $\mathcal{K}$, $P_{\mathcal{H}}$ ortogonalna projekcija na $\mathcal{H}$
    in $T = P_{\mathcal{H}} S\big|_{\mathcal{H}}$.
    Pravimo tudi, da je $T$ \emph{kompresija} $S$.
    Ker je $\mathcal{K} = \mathcal{H} \oplus \mathcal{H}^\perp$, lahko operator $S$ zapišemo v matrični obliki
    $$S = \begin{bmatrix}
        T & *\\
        * & *
    \end{bmatrix}.$$    
\end{definicija}


Pogosto lahko za operator $T$ najdemo dilatacijo, ki ima določene dodatne lastnosti.
To ilustriramo v naslednjih dveh zgledih.


\begin{zgled}\label{zgl:1.1}
    Dokažimo, da ima vsaka izometrija unitarno dilatacijo.
    Naj bo $V \in \bh$ izometrija in $P := I - V V^*$
    projekcija na $(\im V)^\perp$. Zapišimo $\mathcal{K} := \mathcal{H} \oplus \mathcal{H}$ ter 
    definirajmo $U \in \bk$ v matrični obliki kot 
    $$U := \begin{bmatrix}
        V & P\\
        0 & V^*
    \end{bmatrix}.$$
    Očitno je $U$ dilatacija $V$. Dokazati moramo še, da je tudi unitarna. Ker velja
    \begin{align*}
        U^* U &= \begin{bmatrix}
            V^* & 0\\
            P & V
        \end{bmatrix} \begin{bmatrix}
            V & P\\
            0 & V^*
        \end{bmatrix}\\
        &= \begin{bmatrix}
            V^* V &  V^* P\\
            PV & P^2 + V V^*
        \end{bmatrix}\\
        &= \begin{bmatrix}
            I & 0\\
            0 & I
        \end{bmatrix}
    \end{align*}
    in 
    \begin{align*}
         U U^* &= \begin{bmatrix}
            V & P\\
            0 & V^*
        \end{bmatrix} \begin{bmatrix}
            V^* & 0\\
            P & V
        \end{bmatrix} \\
        &= \begin{bmatrix}
             P^2 + V V^* &  P V\\
            V^* P &V^* V 
        \end{bmatrix}\\
        &= \begin{bmatrix}
            I & 0\\
            0 & I
        \end{bmatrix},
    \end{align*}
    je $U$ res unitarna dilatacija $V$. Še več, v tem primeru je 
    $U$ \emph{potenčna dilatacija} $V$. To pomeni, da je za vsak $n \in \N$ operator
    $$U^n = \begin{bmatrix}
        V^n & *\\
        0 & (V^*) ^n
    \end{bmatrix}$$
    dilatacija za $V^n$.
\end{zgled}

\begin{zgled}\label{zgl:1.2}
    Dokažemo lahko tudi, da ima vsaka skrčitev izometrično dilatacijo.
    Vzemimo $T \in \bh$, tako da je $\| T\| \leq 1$. Definirajmo 
    $$D_T := (I - T^* T)^{\frac{1}{2}} \in \bh,$$tako da za vsak $x \in \mathcal{H}$ velja
    \begin{align*}
        \| T x \|^2 + \| D_T x\|^2 &= \langle Tx, Tx \rangle + \langle D_T x, D_T x \rangle\\
        &= \langle T^* T x, x\rangle + \langle D_T ^2 x, x\rangle\\
        &= \langle T^* T x, x\rangle + \langle (I - T^* T) x, x\rangle\\
        &= \langle x, x \rangle = \| x\|^2.
    \end{align*}
    Nato definirajmo Hilbertov prostor $\mathcal{K} := \bigoplus_{n \in \N} \mathcal{H}$, ki je prostor zaporedij
    $$\{(x_1, x_2, \dots)\ |\ x_n \in \mathcal{H},\ \sum_{n = 1} ^\infty \| x_n\|^2 < \infty\}$$
    s skalarnim produktom
    $$\langle (x_1, x_2, \dots), (y_1, y_2, \dots) \rangle := \sum_{n = 1} ^\infty \langle x_i, y_i \rangle.$$
    Nazadnje definiramo še operator 
    $$V: \mathcal{K} \to \mathcal{K},\quad (x_1, x_2,x_3, \dots) \mapsto (T x_1, D_T x_2, x_3, \dots),$$
    ki je izometrija, saj velja
    \begin{align*}
        \| V(x_1, x_2, x_3, \dots)\|^2 &= \| T x_1\|^2 + \| D_T x_2\|^2 + \sum_{n = 3} ^\infty \| x_n\|^2 \\
        &= \sum_{n = 1} ^\infty \| x_n\|^2 \\
        &= \| (x_1, x_2, x_3, \dots )\|^2.
    \end{align*}
    Če identificiramo $\mathcal{H}$ s $\mathcal{H} \oplus 0 \oplus 0 \oplus \cdots \subseteq \mathcal{K}$, potem velja 
    $$T^n = P_{\mathcal{H}} V^n \big|_{\mathcal{H}}$$
    za vse $n \in \N$.
\end{zgled}

Kot posledico zgledov \ref{zgl:1.1} in \ref{zgl:1.2} dobimo naslednji izrek.

\begin{izrek}[Sz.-Nagy]
    Naj bo $T \in \bh$ skrčitev. Potem obstaja Hilbertov prostor $\mathcal{K} \supseteq \mathcal{H}$ 
    in unitaren operator $U \in \bk$, tako da velja
    $$T^n = P_{\mathcal{H}} U^n \big|_{\mathcal{H}}$$
    za vse $n \in \N$.
\end{izrek}

Sz.-Nagyjev izrek nam omogoča, da nekatere trditve o skrčitvah reduciramo na trditve o unitarnih operatorjih.
Za ilustracijo tega argumenta dokažimo naslednjo trditev.

\begin{posledica}[von Neumannova neenakost]\label{pos:1.1}
    Naj bo $T \in \bh$ skrčitev $p \in \C [z]$ kompleksen polinom. Potem imamo
    $$\| p(T)\| \leq \sup \{|p(z)|\ |\ z \in \mathbb{T}\}.$$
\end{posledica}

\begin{dokaz}
    Naj bo $U$ potenčna dilatacija $T$, torej
    $T^n = P_{\mathcal{H}} U^n \big|_{\mathcal{H}}$
    za vse $n \in \N \cup\{0\}$. Potem velja $p(T) = P_{\mathcal{H}} p(U)\big|_{\mathcal{H}}$
    in posledično $\| p(T)\| \leq \| p(U)\|$. Ker je $U$ normalna, imamo po spektralnem izreku
    $$\| p(U) \| = \sup \{|p(\lambda)|\ |\ \lambda \in \sigma(U)\}.$$
    Nato upoštevamo še, da je $U$ unitarna in zanjo velja $\sigma(U) \subseteq \mathbb{T}$, torej
    $$\sup \{|p(\lambda)|\ |\ \lambda \in \sigma(U)\} \leq \sup \{|p(\lambda)|\ |\ \lambda \in \mathbb{T}\}.$$
    Vse to skupaj implicira
    \begin{equation*}
        \| p(T)\| \leq \| p(U)\| = \sup \{|p(\lambda)|\ |\ \lambda \in \sigma(U)\} \leq \sup \{|p(\lambda)|\ |\ \lambda \in \mathbb{T}\}. \qedhere
    \end{equation*} 
\end{dokaz}

\begin{zgled}\label{zgl:1.7}
    Naj bo $T \in \bh$ skrčitev in
    $$\mathcal{M} := \C [z] \subseteq \C (T)$$
    operatorski sistem. Potem je preslikava 
    $$\varphi: \mathcal{M} \to \bh,\quad p \mapsto p(T)$$
    enotska skrčitev po von Neumannovi neenakosti. To pomeni, da 
    jo lahko razširimo do pozitivne preslikave $$\tilde{\varphi}: \mathcal{M} + \mathcal{M}^* \to \bh,\quad p + q^* \mapsto p(T) + q(T)^*.$$
    Po Stone-Weierstrassovem izreku je
    \begin{equation*}
        \mathcal{M} + \mathcal{M}^* = \{p (z) + \overline{q ({z})}\ |\ \text{$p, q$ polinoma}\}. 
    \end{equation*}
    gosta podmnožica v $C(\T)$, zato lahko $\tilde{\varphi}$ razširimo do pozitivne preslikave $\psi: C(\T) \to \bh$.
    Po trditvi \ref{trd:1.2} sledi, da je tudi povsem pozitivna.
\end{zgled}

V prejšnjem zgledu smo dokazali, da za vsako skrčitev $T \in \bh$
obstaja enolična povsem pozitivna preslikava $\psi: C(\T) \to \bh$,
tako da velja $\psi (p) = p(T)$ za poljuben polinom $p \in C(\T)$.
Takoj lahko navedemo dve posledici tega dejstva.

\begin{posledica}[Russo-Dye]\label{pos:russo-dye}
    Če sta $\mathfrak{A}, \mathfrak{B}$ $C^*$-algebri in 
    $\varphi: \mathfrak{A} \to \mathfrak{B}$ pozitivna preslikava, 
    potem je $\| \varphi \| = \| \varphi(1)\|$.
\end{posledica}

\begin{dokaz}
    Brez škode za splošnost predpostavimo, da je $\mathfrak{A} \subseteq \bh$.
    Vzemimo poljuben $a \in \mathfrak{A}$ z $\| a \|\leq 1$. 
    Po prejšnji opombi obstaja povsem pozitivna preslikava $\psi: C(\T) \to \bh$,
    tako da je $\psi(z) = a$. Po konstrukciji je 
    $$\psi (C(\T)) = \overline{\{p(a) |\ p \in \C[z]\}} \subseteq \mathfrak{A}.$$ 
    Posledično je preslikava
    $$\varphi \circ \psi: C(\T) \to \mathfrak{B}$$ pozitivna, 
    zato je po trditvi \ref{trd:1.2} povsem pozitivna 
    Od tod sledi 
    $$\| \varphi(a)\| = \| (\varphi \circ \psi) (z)\| \leq \| (\varphi \circ \psi) (1) \| = \| \varphi (1)\|$$
    in posledično $\| \varphi \| = \| \varphi(1)\|$.
\end{dokaz}

\begin{posledica}[Matrična von Neumannova neenakost]
    Naj bo $T \in \bh$ skrčitev in $p = [p_{ij}]_{i, j} \in M_n(\C [z])$ kompleksna matrični polinom. 
    Potem velja
    $$\| p(T)\| = \| [p_{ij} (T)]_{i, j}\| \leq \sup \{\|[p_{ij}(z)]_{i, j}\|\ |\ z \in \mathbb{T}\}.$$
\end{posledica}

\begin{proof}
    Ker je $\psi: C(\T) \to \bh$ enotska povsem pozitivna preslikava, velja $\| \psi \|_{\po} = 1$.
    V posebnem velja $\| \psi^{(n)} \| =1$ za vsak $n \in \N$. Potem za $[p_{ij}]_{i, j} \in M_n (\C [z])$ velja
    \begin{equation*}
        \| [p_{ij} (T)]_{i, j}\| = \| \psi^{(n)} ([p_{ij} ]_{i, j}) \| \leq \| [p_{ij} ]_{i, j} \|. \qedhere
    \end{equation*}
\end{proof}

Vrnimo se nazaj k dilatacijam.
Na tej točki še ni popolnoma očitno, kaj imajo te skupnega s povsem pozitivnimi preslikavami, na kar smo namignili 
v začetku tega poglavja. Vemo pa, da je za povsem pozitivno preslikavo $\psi: \mathfrak{A} \to \bh$ in $V \in \mathcal{B}(\mathcal{K}, \mathcal{H})$ tudi 
\begin{equation}
    \mathfrak{A} \to  \bk,\quad A \mapsto V^* \psi (a) V \label{eq:1.1}
\end{equation}
povsem pozitivna preslikava. To pomeni, da sta dilatacija in kompresija povsem pozitivni preslikavi na $\bh$.

\begin{zgled}
    Denimo, da sta $\mathcal{H} \subseteq \mathcal{K}$ Hilbertova prostora.
    Potem sta po \eqref{eq:1.1} tudi preslikavi
    $$
        \bk \to \bh,
        \quad
        A \mapsto P_{\mathcal{H}}A|_{\mathcal{H}},
    $$
    in
    $$
        \bh \to \bk,
        \quad
        A \mapsto
        \begin{bmatrix}
            A & 0\\
            0 & 0
        \end{bmatrix}
    $$
    povsem pozitivni. 
    
    Ker je kompozitum povsem pozitivnih preslikav
    ponovno povsem pozitiven, je za vsako povsem pozitivno preslikavo
    $\varphi : \mathfrak{A} \to \bh$
    tudi preslikava
    $$
        \psi : \mathfrak{A} \to \bk,
        \quad
        a \mapsto
        \begin{bmatrix}
            \varphi(a) & 0\\
            0 & 0
        \end{bmatrix}
    $$
    povsem pozitivna.

    Prav tako je za vsako povsem pozitivno preslikavo
    $\psi : \mathfrak{A} \to \bk$
    tudi preslikava
    $$
        \varphi : \mathfrak{A} \to \bh,
        \quad
        \varphi(a)=P_{\mathcal{H}}\psi(a)|_{\mathcal{H}}
    $$
    povsem pozitivna. 
\end{zgled}

Stinespringov izrek nam pove, da je 
vsaka povsem pozitivna preslikava 
$\varphi : \mathcal{A} \to \bh$
oblike \eqref{eq:1.1}. 
Pravzaprav velja še več.

\begin{izrek}[Stinespring]\label{izr:1.1}
    Naj bo $\mathfrak{A}$ $C^*$-algebra in $\varphi: \mathfrak{A} \to \bh$ povsem pozitivna preslikava.
    Potem obstaja Hilbertov prostor $\mathcal{K}$, omejen operator $V \in \mathcal{B} (\mathcal{H}, \mathcal{K})$
    in $*$-homomorfizem 
    $\pi: \mathfrak{A} \to \bk$, tako da je
    $$\varphi(a) = V^* \pi(a) V,\quad \forall a \in \mathfrak{A}.$$
    Če je $\varphi$ enotska preslikava, potem je $V$ izometrija.
\end{izrek}


Za trenutek se ustavimo na primeru, ko je $\varphi: \mathfrak{A} \to \bh$ enotska povsem pozitivna preslikava.
Tedaj je $V$ izometrija, zato lahko
Hilbertov prostor $\mathcal{H}$ identificiramo s podprostorom
$V\mathcal{H} \leq \mathcal{K}$.
S to identifikacijo operator $V^*$ sovpada s projekcijo
$P_{\mathcal{H}}$ iz $\mathcal{K}$ na $\mathcal{H}$, zato velja
$$
    \varphi(a)=P_{\mathcal{H}}\pi(a)|_{\mathcal{H}},\quad \forall a \in \mathfrak{A}.
$$
Stinespringov izrek nam torej pove, da lahko vsako enotsko povsem pozitivno preslikavo $\varphi: \mathfrak{A} \to \bh$ dilatiramo do
$*$-homomorfizma.

\begin{dokaz}[Dokaz izreka \ref{izr:1.1}]
    Oglejmo si algebraični tenzorski produkt vektorskih prostorov $\mathfrak{A} \otimes \mathcal{H}$.
    Na njem definiramo seskvilinearno formo 
    $$\langle a \otimes x, b \otimes y \rangle = \langle \varphi(b^* a) x, y\rangle,$$
    ki jo razširimo linearno. Za vsak $u = \sum_{j = 1} ^n a_j \otimes x_j \in \mathfrak{A} \otimes \mathcal{H}$ imamo
    \begin{align*}
        \langle u, u \rangle &= \sum_{i, j} ^n \langle \varphi(a_i ^* a_j)x_j,x_i \rangle\\
        &= \left\langle \varphi^{(n)} \left(\begin{bmatrix}
            a_1 \\ \vdots \\ a_n
        \end{bmatrix}^* \begin{bmatrix}
            a_1 \\ \vdots \\ a_n
        \end{bmatrix}\right) \begin{bmatrix}
           x_1 \\ \vdots \\x_n
        \end{bmatrix}, \begin{bmatrix}
           x_1 \\ \vdots \\x_n
        \end{bmatrix} \right\rangle \geq 0,
    \end{align*}
    kar pomeni, da je $\langle \cdot, \cdot \rangle$ polskalarni produkt na $\mathfrak{A} \otimes \mathcal{H}$. 

    Vsakemu elementu $a \in \mathfrak{A}$ lahko priredimo preslikavo
    $$\pi_0 (a): \mathfrak{A} \otimes \mathcal{H} \to \mathfrak{A} \otimes \mathcal{H},\quad \sum_{j = 1} ^n a_j \otimes x_j \mapsto \sum_{j = 1} ^n a a_j \otimes x_j.$$
    Sedaj imamo za $u = \sum_{j = 1} ^n a_j \otimes x_j$ in $v = \sum_{i = 1} ^m b_i \otimes y_i$
    \begin{align}
        \langle u, \pi_0 (a) v \rangle &= \langle \sum_{j = 1} ^n a_j \otimes x_j, \sum_{i = 1} ^m a b_i \otimes y_i \rangle \nonumber\\ 
        &= \sum_{i, j} \langle \varphi((a b_i)^* a_j)x_j, y_i \rangle \nonumber\\
        &= \sum_{i, j} \langle \varphi (b_i ^* a^* a_j ) x_j, y_i \rangle \nonumber\\
        &= \langle \pi_0 (a^*) u, v \rangle. \label{eq:1.4}
    \end{align}

    Definirajmo $$\mathcal{N} := \{u \in \mathfrak{A} \otimes \mathcal{H}\ |\ \langle u, u \rangle = 0\}.$$
    Po Cauchy-Schwarzevi neenakosti je $\mathcal{N}$ podprostor v $\mathfrak{A} \otimes \mathcal{H}$ in $\langle \cdot, \cdot \rangle$
    inducira skalarni produkt na $\quot{\mathfrak{A} \otimes \mathcal{H}}{\mathcal{N}}$. 
    Slednji prostor nato dopolnimo do Hilbertovega prostora $\mathcal{K}$.

    Vzemimo poljuben $u \in \mathfrak{A} \otimes \mathcal{H}$ in definirajmo $$f(a) := \langle \pi_0 (a) u, u \rangle.$$ 
    Potem je po \eqref{eq:1.4} $f$ pozitiven linearen funkcional na $\mathfrak{A}$, saj je
    \begin{align*}
        f(a^* a) &= \langle \pi_0 (a^* a) u, u \rangle\\
        &= \langle \pi_0 (a^*) \pi_0 (a) u, u\rangle \\
        &= \langle \pi_0 (a) u, \pi_0 (a) u\rangle \geq 0.
    \end{align*}
    Od tod sledi
    $$\langle \pi_0 (a) u, \pi_0 (a) u\rangle = f(a^* a) \leq \| a^* a\| \cdot f(1) = \| a\|^2 \cdot \langle u, u \rangle.$$
    Posledično je $\pi_0 (a) (\mathcal{N}) \subseteq \mathcal{N}$, torej je $\mathcal{N}$ invarianten podprostor za $\pi_0 (a)$ v $\mathfrak{A} \otimes \mathcal{H}$, zaradi česar 
    $\pi_0 (a)$ inducira omejen operator na $\quot{\mathfrak{A} \otimes \mathcal{H}}{\mathcal{N}}$,
    ki ga nato razširimo do omejenega operatorja na $\mathcal{K}$. 
    Zato $\pi_0$ inducira preslikavo (ki je hkrati $*$-homomorfizem) 
    $$\pi:\mathfrak{A} \to \bk,\quad \pi (a) (u + \mathcal{N}) = \pi_0 (a)u + \mathcal{N}.$$

    Nazadnje definiramo še $$V: \mathcal{H} \to \mathcal{K},\quad x \mapsto 1 \otimes x + \mathcal{N}.$$
    Ker je $$\| V x \|^2 = \langle 1 \otimes x, 1 \otimes x \rangle_{\mathcal{K}} = \langle \varphi (1) x, x \rangle_{\mathcal{H}} \leq \| \varphi (1)\| \cdot \| x \|^2,$$
    je $V$ omejen operator. Za poljubna $x, y \in \mathcal{H}$ in $a \in \mathfrak{A}$ imamo
    \begin{align*}
        \langle V x, a \otimes y + \mathcal{N} \rangle &= \langle 1 \otimes x + \mathcal{N}, a \otimes y + \mathcal{N}\rangle\\
        &= \langle \varphi (a^*) x, y \rangle\\
        &= \langle x, \varphi (a^*)^* y \rangle\\
        &= \langle x, \varphi(a) y \rangle,
    \end{align*}
    torej je $V^* (a \otimes y + \mathcal{N}) = \varphi(a) y$. Sedaj je
    $$V^* \pi(a) V y = V^* \pi (a) (1 \otimes y + \mathcal{N}) = V^* (a \otimes y + \mathcal{N}) = \varphi(a) y.$$

    Če je $\varphi$ enotska preslikava, potem za poljuben vektor $y \in \mathcal{H}$ velja
    $$V^* V y = V^* \pi(1) V y = \varphi(1) y = y,$$
    torej je $V^* V = I$ in je $V$ izometrija.
\end{dokaz}

\begin{opomba}
    Stinespringov izrek je posplošitev izreka GNS o upodobitvah. 
    Če namreč vzamemo $\mathcal{H} = \C$, potem je $\mathcal{B}(\C) = \C$ in $\mathfrak{A} \otimes \C = \mathfrak{A}$, tako da je izometrija 
    $V: \C \to \mathcal{K}$ določena z $V(1) = x$. Tako imamo 
    $$\varphi(a) = \varphi(a) (1) \cdot 1 = V^* \pi(a) V(1) \cdot 1 = \langle \pi (a) V(1), V(1) \rangle_{\mathcal{K}}$$
\end{opomba}

Trojici $(\pi, V, \mathcal{K})$ iz Stinespringovega izreka pravimo \emph{Stinespringova upodobitev}
za $\varphi$. Naravno se je vprašati, ali je Stinespringova upodobitev enolična. 
Če že imamo eno Stinespringovo upodobitev $(\pi, V, \mathcal{K})$ za $\varphi$, lahko definiramo nov prostor $\mathcal{K}_1 := \overline{\pi(\mathfrak{A}) V \mathcal{H}} \leq \mathcal{K}$, ki reducira $\pi (\mathfrak{A})$.
To nam omogoča, da definiramo novo upodobitev $\pi_1: \mathfrak{A} \to \mathcal{B}(\mathcal{K}_1)$, za katero prav tako velja $$\varphi(a) = V^* \pi_1 (a) V,\quad \forall a \in \mathfrak{A}$$
in $(\pi_1, V, \mathcal{K}_1)$ je še ena Stinespringova upodobitev za $\varphi$. To je motivacija za naslednjo definicijo.

\begin{definicija}
    Stinespringova upodobitev $(\pi, V, \mathcal{K})$ povsem pozitivne preslikave $\varphi: \mathfrak{A} \to \bh$ je \emph{minimalna}, če je $\overline{\pi(\mathfrak{A}) V \mathcal{H}} = \mathcal{K}$.
\end{definicija}

\begin{trditev}
    Naj bo $\mathfrak{A}$ $C^*$-algebra in $\varphi: \mathfrak{A} \to \bh$ povsem pozitivna preslikava. 
    Če sta $(\pi_1, V_1, \mathcal{K}_1), (\pi_2, V_2, \mathcal{K}_2)$ 
    minimalni Stinespringovi upodobitvi, potem obstaja unitaren operator $U: \mathcal{K}_1 \to \mathcal{K}_2$, tako da velja
    $U V_1 = V_2$ in $U \pi_1 U^* = \pi_2$.
\end{trditev}

\begin{dokaz}
    Definiramo preslikavo
    $$\widetilde{U} : \pi_1 (\mathfrak{A}) V_1 \mathcal{H} \to \pi_2 (\mathfrak{A}) V_2 \mathcal{H},\quad \sum_i \pi_1 (a_i) V_1 x_i \mapsto \sum_i \pi_2 (a_i) V_2 x_i.$$
    Takoj lahko opazimo, da je $\widetilde{U}$ surjekcija.

    Nadalje dokažemo, da je $\widetilde{U}$ izometrija:
    \begin{align*}
        \left\| \sum_i \pi_1 (a_i) V_1 x_i \right\|^2 &= \left\langle \sum_i \pi_1 (a_i) V_1 x_i, \sum_j \pi_1 (a_j) V_1 x_j \right\rangle\\
        &= \sum_i \sum_j \langle \pi_1 (a_i) V_1 x_i, \pi_1 (a_j) V_1 x_j\rangle\\
        &= \sum_i \sum_j \langle V_1^* \pi_1 ^* (a_j) \pi_1 (a_i) V_1 x_i,  x_j\rangle\\
        &= \sum_i \sum_j \langle V_1^* \pi_1 (a_j^* a_i) V_1 x_i,  x_j\rangle\\
        &= \sum_i \sum_j \langle \varphi(a_j^* a_i) x_i,  x_j\rangle\\
        &= \sum_i \sum_j \langle V_2^* \pi_2 (a_j^* a_i) V_2 x_i,  x_j\rangle\\
        &= \left\langle \sum_i \pi_2 (a_i) V_2 x_i, \sum_j \pi_2 (a_j) V_2 x_j \right\rangle\\
        &= \left\| \sum_i \pi_2 (a_i) V_2 x_i \right\|^2.
    \end{align*}
    Torej lahko razširimo $\widetilde{U}$ do izometrije $U: \mathcal{K}_1 \to \mathcal{K_2}$
    Ker je $U(\mathcal{K}_1)$ gosta v $\mathcal{K}_2$, je $U$ surjektivna.

    Enakost $U V_1 = V_2$ sledi direktno iz definicije $U$.
    Do dokažemo $U \pi_1 = \pi_2 U$, opazimo, da za vsak $a \in \mathfrak{A}$ velja 
    \begin{align*}
        \pi_2 (a) U \big|_{\pi_1 (\mathfrak{A}) V_1 \mathcal{H}}
        &=  \pi_2 (a) \widetilde{U} \big|_{\pi_1 (\mathfrak{A}) V_1 \mathcal{H}}\\
        &= \widetilde{U} \pi_1 (a) \big|_{\pi_1 (\mathfrak{A}) V_1 \mathcal{H}}\\
        &= U \pi_1 (a) \big|_{\pi_1 (\mathfrak{A}) V_1 \mathcal{H}},
    \end{align*}
    torej $\pi_2 (a) U = U \pi_1 (a)$. 
\end{dokaz}

\begin{opomba}
    V dokazu izreka \ref{izr:1.1} je množica
    $$\pi(\mathfrak{A}) V \mathcal{H} = \quot{\mathfrak{A} \otimes \mathcal{H}}{\mathcal{N}}$$
    gosta v svoji napolnitvi $\mathcal{K}$, torej v dokazu konstruiramo minimalno Stinespringovo upodobitev.
\end{opomba}


\begin{posledica}[Sz.-Nagyjeva minimalna unitarna dilatacija] \label{pos:1.4}
    Naj bo $T \in \bh$ skrčitev. Potem obstaja Hilbertov prostor $\mathcal{K} \supseteq \mathcal{H}$
    in unitaren operator $U \in \bk$, tako da je $\mathcal{K}$ najmanjši reducirajoč prostor za $U$, da velja 
    $$T^n = P_{\mathcal{H}} U^n \big|_{\mathcal{H}},\quad \forall n \in \N.$$

    Če je $(U', \mathcal{K}')$ še en takšen par,
    potem obstaja unitaren operator $V: \mathcal{K} \to \mathcal{K}'$, tako da je $V\big|_{\mathcal{H}} = \id_{\mathcal{H}}$
    in $V U V^* = U'$. 
\end{posledica}

\begin{dokaz}
    Stinespringov izrek uporabimo na pozitivni preslikavi $\psi: C(\T) \to \bh$ iz zgleda \ref{zgl:1.7},
    ki preslika poljuben polinom $p \in \C[z]$ v $p(T) \in \bh$. 
    Potem obstaja minimalna Stinespringova upodobitev $(\pi, V, \mathcal{K})$,
    pri čemer je $V$ izometrija, saj je $\psi$ enotska preslikava. 
    Posledično lahko identificiramo $\mathcal{H} \equiv V \mathcal{H}$.

    Sedaj za poljuben $n \in \N$ velja
    $$T^n = \psi(z^n) = P_{\mathcal{H}} \pi(z) ^n \big|_{\mathcal{H}} = P_{\mathcal{H}} U ^n \big|_{\mathcal{H}},$$
    kjer je $U := \pi(z) \in \bk$ unitaren operator.
    Po zveznem funkcijskem računu je $\pi(f) = f(U)$ za poljubno funkcijo $f \in C(\T)$.
    Ker je $(\pi, V, \mathcal{K})$ minimalna Stinespringova upodobitev, je $\overline{\pi(C(\T)) \mathcal{H}} = \mathcal{K}$.
    Ker pa je $\pi(C(\T)) = C^*(U)$, je to ekvivalentno $$\overline{\{U^n h\ |\ n \in \Z,\ h \in \mathcal{H}\}} = \mathcal{K}.$$
    Slednje pa je ekvivalentno temu, da je $\mathcal{K}$ najmanjši reducirajoč prostor za $U$.
    Zadnja trditev v izreku nato sledi direktno iz unitarne ekvivalence minimalnih Stinespringovih upodobitev.
\end{dokaz}


Razdelek zaključimo s posledico Stinespringovega izreka, ki jo bomo potrebovali kasneje.
Gre za `nekomutativno posplošitev' Radon-Nikodymovega izreka.
Najprej na množici $\PP(\mathfrak{A}, \mathfrak{B})$ definiramo delno urejenost  
$$\varphi \leq \psi \quad \Leftrightarrow \quad \psi - \varphi \in \PP(\mathfrak{A}, \mathfrak{B}).$$
Nadalje za poljubno podmnožico $C^*$-algebre $\mathcal{S} \subseteq \mathfrak{A}$ definiramo njen \emph{komutant} kot
$$\mathcal{S}' := \{a \in A\ |\ as = sa,\ \forall s \in \mathcal{S}\}.$$
Sedaj velja naslednje.

\begin{trditev} \label{trd:1.4}
    Naj bosta $\varphi, \psi \in \PP(\mathfrak{A},\bh)$ in
    $(\pi, V, \mathcal{K})$ minimalna Stinespringova upodobitev $\psi$.
    Potem je $\psi \geq \varphi$ natanko tedaj, ko obstaja enoličen $T \in \pi(\mathfrak{A})'$,
    tako da je
    $$\varphi(a) = V^* T \pi(a) V$$ in $0 \leq T \leq 1$.
\end{trditev}

\begin{dokaz}
    Denimo, da je $\varphi = V^* T \pi(a) V$, pri čemer velja $0 \leq T \leq 1$ in $T \in \pi(\mathfrak{A})'$.
    Ker je $1 - T \geq 0$, lahko definiramo $S := (1 - T)^{\frac{1}{2}}$. 
    Opazimo, da velja $1 - T \in \pi(\mathfrak{A})'$, zato imamo po zveznem funkcijskem računu $S \in \pi(\mathfrak{A})'$.
    Od tod pa sledi, da je $$\psi - \varphi = V^* (1 - T) \pi(a) V = V^* S^* \pi(a) S V$$
    povsem pozitivna.
    
    Za obratno implikacijo predpostavimo, da je $\psi - \varphi$ povsem pozitivna.
    Ker so vse minimalne Stinespringove dilatacije $\psi$ unitarno ekvivalentne, lahko brez škode za splošnost vzamemo trojico $(\pi, V, \mathcal{K})$,
    ki smo jo konstruirali v dokazu izreka \ref{izr:1.1}.
    V istem dokazu smo vpeljali polskalarna produkta 
    $$\langle a \otimes x, b \otimes y \rangle_{\varphi} := \langle \varphi(b^* a) x, y\rangle,\quad \langle a \otimes x, b \otimes y \rangle_{\psi} := \langle \psi(b^* a) x, y\rangle$$
    na prostoru $\mathfrak{A} \otimes \mathcal{H}$.
    Ker je $\psi - \varphi$ povsem pozitivna, sledi $$\langle u, u \rangle_{\varphi} \leq \langle u, u \rangle_{\psi}$$
    za poljuben $u \in \mathfrak{A} \otimes \mathcal{H}$. Po Cauchy-Schwarzevi neenakosti imamo za poljubna $u, v \in \mathcal{H}$
    \begin{equation}
        \langle u, v\rangle_{\varphi} \leq \langle u, u\rangle_{\varphi} \langle v, v \rangle_{\varphi} \leq \langle u, u\rangle_{\psi} \langle v, v \rangle_{\psi}. \label{eq:1.5}
    \end{equation} 
    Označimo $\mathcal{N} = \{u \in \mathfrak{A} \otimes \mathcal{H}\ |\ \langle u, u \rangle_{\psi}\}$, tako da $\langle \cdot, \cdot \rangle_{\psi}$ tvori skalarni produkt na prostoru $\quot{\mathfrak{A} \otimes \mathcal{H}}{\mathcal{N}}.$
    Slednjega nato napolnimo do Hilbertovega prostora $\mathcal{K}$.

    Po neenakosti \eqref{eq:1.5} $\langle \cdot, \cdot \rangle_\varphi$ inducira omejen polskalarni produkt na $\mathcal{K}$.
    Po Rieszovem izreku obstaja operator $T \in \bk$, tako da je
    $$\langle u, v \rangle_{\varphi} = \langle T u, v \rangle_{\psi},\quad \forall u, v \in \mathcal{K}.$$
    Potem za vsak $a \in \mathfrak{A}$ in $x, y \in \mathcal{H}$ velja
    \begin{align*}
        \langle (V^* T \pi(a) V) x, y\rangle &= \langle T \pi(a) V x, V y\rangle_{\psi}\\
        &= \langle T \pi(a) (1 \otimes x + \mathcal{N}), (1 \otimes y + \mathcal{N}) \rangle _{\psi}\\
        &= \langle T (a \otimes x + \mathcal{N}), (1 \otimes y + \mathcal{N}) \rangle _{\psi}\\
        &= \langle (a \otimes x + \mathcal{N}), (1 \otimes y + \mathcal{N}) \rangle _{\varphi}\\
        &= \langle \varphi (a) x, y \rangle,
    \end{align*}
    od koder sledi $\varphi(a) = V^* T \pi(a) V$. Prav tako za vsak $u \in \mathcal{K}$ velja
    $$0 \leq \langle T u, u \rangle_{\psi} = \langle u, u \rangle_{\varphi} \leq \langle u, u \rangle_{\psi},$$
    torej $0 \leq T \leq 1$. 
    
    Da dokažemo $T \in \pi(\mathfrak{A})'$, vzamemo poljubne $a, b, c \in \mathfrak{A}$ ter $x, y \in \mathcal{H}$.
    Potem \begin{align*}
        \langle T \pi(a) (b \otimes x + \mathcal{N}), (c \otimes y + \mathcal{N}) \rangle_\psi &= \langle T (ab \otimes x + \mathcal{N}), (c \otimes y + \mathcal{N}) \rangle_\psi\\
        &= \langle (ab \otimes x + \mathcal{N}), (c \otimes y + \mathcal{N}) \rangle_\varphi\\
        &= \langle \varphi (c^* ab)x, y\rangle\\
        &= \langle \varphi ((a^* c)^* b)x, y \rangle\\
        &= \langle (b \otimes x + \mathcal{N}), (a^* c \otimes y + \mathcal{N}) \rangle_\varphi\\
        &= \langle T (b \otimes x + \mathcal{N}), \pi(a^*) (c \otimes y + \mathcal{N}) \rangle_{\psi}\\
        &= \langle\pi(a) T  (b \otimes x + \mathcal{N}), (c \otimes y + \mathcal{N}) \rangle_\psi     
    \end{align*}
    torej $\langle T \pi(a) u, v \rangle_{\psi} = \langle \pi(a) T u, v \rangle_{\psi}$ za vse $u, v \in \mathcal{K}$. Od tod dobimo $T \pi(a) = \pi(a) T$.

    Za dokaz enoličnosti vzemimo nek $T' \in \pi(\mathfrak{A})'$, tako da je $0 \leq T' \leq 1$ in $\varphi = V^* T' \pi V$.
    Ponovno vzemimo poljubna $a, b \in \mathfrak{A}$ in $x, y \in \mathcal{H}$. Potem je 
    \begin{align*}
        \langle T (a \otimes x + \mathcal{N}), (b \otimes y + \mathcal{N}) \rangle _{\psi} &= \langle (V^* T \pi(b^* a) V) x, y\rangle\\
        &= \langle (V^* T' \pi(b^* a) V) x, y\rangle\\
        &= \langle T' (a \otimes x + \mathcal{N}), (b \otimes y + \mathcal{N}) \rangle _{\psi},
    \end{align*}
    zato je $\langle Tu, v \rangle_{\psi} = \langle T' u, v\rangle_{\psi}$ za poljubna $u, v \in \mathcal{K}$.
    Od tod dobimo $T' = T$.
\end{dokaz}



\begin{comment}


Nadaljujemo s posplošitvijo Radon-Nikodymovega izreka za povsem pozitivne preslikave.

\begin{definicija}
    Na množici $\PP(\mathfrak{A}, \mathfrak{B})$ lahko definiramo delno urejenost  
    $$\varphi \leq \psi \quad \Leftrightarrow \quad \psi - \varphi \in \PP(\mathfrak{A}, \mathfrak{B}).$$
\end{definicija}


Za zaključek podrazdelka dokažemo uporabno posledico Stinespringovega izreka, ki klasificira vse povsem pozitivne preslikave med 
prostori končno razsežnih matrik.

\begin{izrek}[Choi]\label{izr:1.2}
    Naj bo $\varphi: M_n (\C) \to M_m (\C)$ povsem pozitivna.
    Potem obstaja $r \leq m \cdot n$, tako da imamo $r$ kompleksnih matrik $V_1, \dots, V_r$ dimenzije $n \times m$,
    ki zadoščajo 
    \begin{equation}
        \varphi(A) = \sum_{k = 1} ^r V_k ^* A V_k,\quad \forall A \in M_n (\C). \label{eq:1.6}
    \end{equation}
\end{izrek}

Za dokaz potrebujemo naslednjo lemo.

\begin{lema}
    Naj bo $\mathfrak{A}$ $C^*$-algebra in $\pi: M_n (\mathfrak{A}) \to \bk$ upodobitev.
    Potem obstaja upodobitev $\rho: \mathfrak{A} \to \bh$, tako da je  
    $\pi$ unitarno ekvivalentna $\rho^{(n)}: M_n (\mathfrak{A}) \to \mathcal{B}(\mathcal{H}^n)$.
\end{lema} 

\begin{proof}
    Za poljuben par indeksov $i, j$ definirajmo $P_{ij} := \pi (E_{ij}) \in \bk$.
    Očitno velja $P_{ij} ^* = P_{ji}$ in 
    \begin{equation}\label{eq:1.7}
       P_{ij} P_{kl} = \delta_{jk} P_{il}. 
    \end{equation}
    V posebnem je 
    $$P_{ii} ^2 = P_{ii} = P_{ii}^*,$$
    torej je $P_{ii}$ ortogonalna projekcija na zaprt (po izreku o zaprtem grafu)
    podprostor $\mathcal{H}_i \subseteq \mathcal{K}$.
    Ker je $$P_{11} + \cdots + P_{nn} = I,\quad P_{ii} P_{jj} = 0,$$
    mora veljati $$\mathcal{H}_1 + \cdots + \mathcal{H}_n = \mathcal{K},\quad \mathcal{H}_i \perp \mathcal{H}_j$$
    in posledično $$\mathcal{H}_1 \oplus \cdots \oplus \mathcal{H}_n = \mathcal{K}.$$
    Za vsak $a \in \mathfrak{A}$ ter $P_{ii}$ velja $P_{ii} \pi(aI) = \pi(aI) P_{ii}$, torej je vsak $\mathcal{H}_i$
    reducirajoč prostor za $\pi(aI)$ in je 
    $$\pi_{ii}: \mathfrak{A} \to \mathcal{B}(\mathcal{H}_i),\quad \pi_{ii} (a) = \pi(aI) \big|_{\mathcal{H}_i}$$
    upodobitev.

    Po enačbi \eqref{eq:1.7} imamo $P_{ij} (\mathcal{H}_k) = \delta_{jk} \mathcal{H}_{i}$, 
    kar pomeni, da lahko $P_{ij}$ zapišemo v matrični obliki na $\mathcal{H}_1 \oplus \cdots \oplus \mathcal{H}_n$ kot
    $$P_{ij} = \begin{bmatrix}
    0 & 0 & 0 \\
    0  & 0 & \underbrace{P_{ij}\big|_{\mathcal{H}_j}}_{(i, j)} \\
     0 & 0 & 0 
 \end{bmatrix}.$$
    Iz \eqref{eq:1.7} prav tako sledi $$P_{ij} ^* P_{ij} = P_{ji} P_{ij} = P_{jj},$$
    od koder dobimo
    $$(P_{ij}\big|_{\mathcal{H}_j})^* (P_{ij}\big|_{\mathcal{H}_j}) = (P_{ji}\big|_{\mathcal{H}_i}) (P_{ij}\big|_{\mathcal{H}_j}) = P_{jj} \big|_{\mathcal{H}_j} = \id_{\mathcal{H}_j}.$$
    To implicira, da je $P_{ij}\big|_{\mathcal{H}_j}: \mathcal{H}_j \to \mathcal{H}_i$ unitaren operator
    in $$P_{ij} \pi_{jj}(a) P_{ji} = \pi_{ii} (a),\quad \forall a \in \mathfrak{A}.$$

    Sedaj definiramo $\mathcal{H} := \mathcal{H}_1$, $\rho := \pi_{11}$ in $U_i := P_{i1} \big|_{\mathcal{H}_1} : \mathcal{H}_1 \to\mathcal{H}_i$, zato je 
    $$U = \oplus_i ^n U_i : \mathcal{H}^n \to \oplus_{i = 1} ^n \mathcal{H}_i = \mathcal{K}$$
    unitaren operator. Ker je
    \begin{align*}
        U^* \pi(a E_{ij}) U &= \begin{bmatrix}
            U_1 ^* & & \\
            & \ddots & \\
            & & U_n ^* 
        \end{bmatrix} \pi_{ii} (a) P_{ij} \begin{bmatrix}
            U_1  & & \\
            & \ddots & \\
            & & U_n  
        \end{bmatrix}\\
        &= \begin{bmatrix}
            P_{11}\big|_{\mathcal{H}_1}  & & \\
            & \ddots & \\
            & & P_{1n} \big|_{\mathcal{H}_n} 
        \end{bmatrix} \begin{bmatrix}
           0 & 0 & 0 \\
           0  & 0 & \pi_{ii} (a) {P_{ij}\big|_{\mathcal{H}_j}} \\
            0 & 0 & 0 
        \end{bmatrix}
        \begin{bmatrix}
            P_{11}\big|_{\mathcal{H}_1}  & & \\
            & \ddots & \\
            & & P_{n1} \big|_{\mathcal{H}_1} 
        \end{bmatrix}\\
        &= \begin{bmatrix}
            0 & 0 & 0 \\
            0  & 0 & (P_{1i} \big|_{\mathcal{H}_i}) \pi_{ii} (a){(P_{ij}\big|_{\mathcal{H}_j})}(P_{j1} \big|_{\mathcal{H}_1}) \\
             0 & 0 & 0 
         \end{bmatrix}\\
         &= \begin{bmatrix}
            0 & 0 & 0 \\
            0  & 0 & (P_{1i} \big|_{\mathcal{H}_i}) \pi_{ii} (a)(P_{i1} \big|_{\mathcal{H}_1}) \\
             0 & 0 & 0 
         \end{bmatrix}\\
         &= \begin{bmatrix}
            0 & 0 & 0 \\
            0  & 0 & \pi(a)  \\
             0 & 0 & 0 
         \end{bmatrix},
    \end{align*}
    sledi $U^* \pi([a_{ij}]_{i, j}) U = [\pi(a_{ij})]_{i, j}$.
\end{proof}

Če v zgornji lemi vzamemo $\mathfrak{A} = \C$, potem za vsak $*$-homomorfizem 
$\pi: M_n (\C) \to \bk$ obstaja Hilbertov prostor $\mathcal{H}$ in unitaren operator 
$U: \mathcal{H}^n \to \mathcal{K}$, tako da za vsako matriko
$A = [a_{ij}]_{i, j} \in M_n (\mathfrak{A})$ velja
$$U^* \pi(A) U = A \otimes \id_{\mathcal{H}} \in M_n (\bh) = \mathcal{B} (\mathcal{H}^n).$$

Denimo, da je $\{v_i\}_{i \in I}$ ortonormalna baza $\mathcal{H}$.
Potem imamo dekompozicijo $$\mathcal{H}^n = \oplus _{i \in I} \mathcal{V}_i,$$
pri čemer je $\mathcal{V}_i = \lin \{u_i\}^n$. Zgornja lema nam potem pove, da je
$$U^* \pi(A) U \big|_{\mathcal{V}_i} = A : \mathcal{V}_i \to \mathcal{V}_i,$$
od koder sledi $U^* \pi(A) U = \bigoplus_{i \in I} A$.

\begin{dokaz}[Dokaz izreka \ref{izr:1.2}]
    Naj bo $\varphi: M_n (\C) \to M_m (\C) = \mathcal{B}(\C^m)$ povsem pozitivna preslikava.
    Po Stinespringovem izreku obstaja upodobitev $\pi: M_n (\C) \to \bk$ in $V : \C^m \to \mathcal{K}$, 
    tako da je
    $$\varphi(A) = V^* \pi(A) V$$ za vse $A \in M_n (\C)$.
    Iz konstrukcije Stinespringovega izreka sledi $$\dim (\mathcal{K}) \leq \dim (M_n (\C) \otimes \C^m) = n^2 m.$$
    Po prejšnjem odstavku lahko vzamemo $\mathcal{K} = \bigoplus_{k = 1} ^r \C ^{n} _k$,
    pri čemer je $r \leq nm$ in $\pi(A) \big|_{\C_k ^n} = A$ za vsak $k \leq r$. V bločno matrični obliki preslikavo $V: \C^m \to \bigoplus_{k = 1} ^r \C ^{n} _k$ sedaj zapišemo
    $$V = \begin{bmatrix}
        V_1 \\ \vdots \\ V_r
    \end{bmatrix},\quad  V_i \in \C^{n \times m}.$$
    Če vse to združimo, dobimo
    \begin{align*}
        \varphi(A) &= V^* \pi(A) V\\
        &= \begin{bmatrix}
            V_1 ^* & \cdots & V_r ^*
        \end{bmatrix} \begin{bmatrix}
            A &  & \\
             & \ddots &  \\
             &  & A
        \end{bmatrix}
        \begin{bmatrix}
            V_1 \\ \vdots \\ V_r
        \end{bmatrix} \\
        &= \sum_{k = 1} ^r V_k ^* A V_k
    \end{align*}
    za vsako matriko $A \in M_n (\C)$.
\end{dokaz}

\end{comment}

\section{Povsem pozitivne preslikave na $M_n (\C)$}\label{sec:7}


V tem oglavju karakteriziramo povsem pozitivne preslikave $\varphi: \mathfrak{A} \to \mathfrak{B}$,
kjer je $\mathfrak{A}$ ali $\mathfrak{B}$ oblike $M_n (\C)$. 
Kot posledica izpeljane teorije bo sledilo, da lahko vsako povsem pozitivno preslikavo $\varphi: \mathcal{S} \to M_n (\C)$ 
na poljubnem operatorskem sistemu $\mathcal{S} \subseteq \mathfrak{A}$ razširimo do povsem pozitivne preslikave na celotni $C^*$-algebri $\mathfrak{A}$.
To je ključni rezultat, ki ga bomo potrebovali za dokaz Arvesonovega razširitvenega izreka v naslednjem poglavju.

Začnimo najprej s karakterizacijo povsem pozitivnih preslikav $\varphi: M_n (\C) \to M_m (\C)$.
Ključen pojem bo \emph{Choijeva matrika} $\varphi^{(n)}([E_{ij}]_{i, j})$, kjer $E_{ij}$ označujejo standardne matrične bazne enote v $M_n (\C)$.
Če definiramo vektor 
$$e = \begin{bmatrix}
    e_1 \\ \vdots \\ e_n
\end{bmatrix} \in \C^{n^2},$$
kjer so $e_i$ standardni bazni vektorji v $\C^n$, potem vidimo, da je 
$$e e^* = [E_{ij}]_{i, j} \in M_n (M_n (\C))$$
pozitivna matrika. 
Od tod sledi, da ima povsem pozitivna preslikava $\varphi: M_n (\C) \to M_m(\C)$ vedno pozitivno Choijevo matriko.
Naslednja trditev, ki je bila dokazana v \cite{kraus1971} in \cite{choi1975}, nam pove, da velja tudi obratno.


\begin{comment}
\begin{opomba}
    V zgornjem dokazu smo dokazali, da za vsako pozitivno matriko $A \in M_{mn} (\C)$ obstajajo matrike 
    $V_1, \dots, V_{mn} \in \C^{n \times m}$, tako da je 
    $$A = \sum_{k = 1} ^{mn} (I_n \otimes V_k)^* [E_{ij}]_{i, j} (I_n \otimes V_k).$$
\end{opomba}
\end{comment}

\begin{izrek}[Choi-Kraus]\label{izr:1.2}
    Za linearno preslikavo $\varphi: M_n (\C) \to M_m (\C)$ so naslednje trditve ekvivalentne:
    \begin{enumerate}
        \item $\varphi$ je povsem pozitivna;
        \item $\varphi$ je $n$-pozitivna;
        \item {Choijeva matrika} $\varphi^{(n)}([E_{ij}]_{i, j})$ je pozitivna;
        \item $\varphi$ je oblike 
        \begin{equation}
            \varphi(A) = \sum_{k = 1} ^r V_k ^* A V_k,\quad \forall A \in M_n (\C). \label{eq:1.6}
        \end{equation}
    \end{enumerate}
\end{izrek}


\begin{dokaz}
    Edina netrivialna implikacija je $(3.) \Rightarrow (4.)$, zato predpostavimo,
    da je Choijeva matrika $\varphi^{(n)}([E_{ij}]_{i, j})$ pozitivna.
    Po lemi \ref{lem:1.2} obstajajo vrstični vektorji $v_1, \dots, v_{mn} \in \C^{mn}$, tako da je 
    $$[\varphi (E_{ij})]_{i, j} = \sum_{k = 1} ^{mn} v_k v_k ^* \in M_n (M_m (\C)) = M_{mn} (\C).$$
    Vsakemu stolpcu
    $$v_k = \begin{bmatrix}
        x_1 ^{(k)}  \\ \vdots \\ x_n ^{(k)}
    \end{bmatrix} \in \C^{mn},\quad x_j ^{(k)} \in \C^{ m},$$
    priredimo matriko 
    $$V_k = \begin{bmatrix}
        {x_1 ^{(k)}} ^* \\
         \vdots \\
         {x_n ^{(k)} }^*
    \end{bmatrix} \in \C^{n \times m}.$$
    Opazimo, da velja
    $$V_k ^* E_{ij} V_k  = (V_k ^* e_i) (e_j^* V_k)  = {x_i ^{(k)}} {x_j ^{(k)}} ^* = v_k v_k ^*,$$
    torej
    $$\varphi^{(n)} ([E_{ij}]_{i, j}) = \sum_{k = 1} ^{mn} (I_n \otimes V_k)^* [E_{ij}]_{i, j} (I_n \otimes V_k).$$ 
    Sedaj za vsako matriko $A = \sum_{i, j = 1} ^n a_{ij} E_{ij} \in M_{n} (\C)$ velja
    \begin{align*}
        \varphi (A) &= \varphi \left(\sum_{i, j = 1} ^n a_{ij} E_{ij}\right) = \sum_{i, j = 1} ^n a_{ij} \varphi \left(E_{ij}\right)\\
        &= \sum_{i, j = 1} ^n a_{ij} \sum_{k = 1}^{mn} V_k ^* E_{ij} V_k = \sum_{k = 1}^{mn} \sum_{i, j = 1} ^n a_{ij} V_k ^* E_{ij} V_k\\
        &= \sum_{k = 1}^{mn} V_k ^* \left( \sum_{i, j = 1} ^n a_{ij} E_{ij} \right) V_k = \sum_{k = 1}^{mn} V_k ^* A V_k. \qedhere
    \end{align*}
\end{dokaz}


Naj bo $\mathfrak{B}$ sedaj poljubna $C^*$-algebra.
S pomočjo Choijeve matrike lahko karakteriziramo povsem pozitivne preslikave $\varphi: M_n (\C) \to \mathfrak{B}$.

\begin{trditev}\label{trd:1.6}
    Za linearno preslikavo $\varphi: M_n (\C) \to \mathfrak{B}$ so naslednje
    trditve ekvivalentne:
    \begin{enumerate}
        \item $\varphi$ je povsem pozitivna;
        \item $\varphi$ je $n$-pozitivna;
        \item Choijeva matrika $[\varphi (E_{ij})]_{i, j} \in M_n (\mathfrak{B})$ je pozitivna.
    \end{enumerate}
\end{trditev}

\begin{dokaz}
    Dovolj je dokazati implikacijo $(3.) \Rightarrow (1.)$.
    Brez škode za splošnost predpostavimo, da je $\mathfrak{B} = \bh$ za nek Hilbertov prostor $\mathcal{H}$.
    Naj bo $[A_{ij}]_{i, j} \in M_m (M_n (\C))$ pozitivna matrika.
    Vzemimo poljubne vektorje $x_1, \dots, x_m \in \mathcal{H}$ in dokazujemo, da je 
    $$\sum_{i, j} ^m \langle \varphi (A_{ij}) x_j, x_i \rangle \geq 0.$$
    Definiramo $$V: \C^m \to \mathcal{H},\quad V e_i = x_i$$ in
    $$\psi : M_n (\C) \to M_m (\C),\quad \psi (A) = V^* \varphi(A) V.$$
    Ker je Choijeva matrika 
    $$\psi^{(n)} ([E_{ij}]_{i, j}) = [V^* \varphi (E_{ij}) V]_{i, j} = (I_n \otimes V)^* [\varphi (E_{ij})]_{i, j} (I_n \otimes V)$$
    po predpostavki pozitivna, je po Choijevem izreku $\psi$ povsem pozitivna preslikava. 
    Od tod pa sledi 
    \begin{equation*}
        \sum_{i, j} ^m \langle \varphi (A_{ij}) x_j, x_i \rangle = \sum_{i, j} ^m \langle V^* \varphi (A_{ij}) V e_j, e_i \rangle = \sum_{i, j} ^m \langle \psi (A_{ij}) e_j, e_i \rangle \geq 0. \qedhere
    \end{equation*}
\end{dokaz}

Naš cilj za preostanek tega poglavja je pridobiti podobno karakterizacijo povsem pozitivnih preslikav $\varphi: \mathfrak{A} \to M_n (\C)$.

Naj bo $\mathcal{M} \subseteq \mathfrak{A}$ operatorski prostor.
Za poljubno linearno preslikavo $\varphi \in \mathcal{L} (\mathcal{M}, M_n (\C))$
lahko definiramo $$s_{\varphi} \in \mathcal{L}(M_n (\mathcal{M}), \C),\quad s_{\varphi} (A) = \frac{1}{n} \sum_{i, j} ^n \varphi(a_{ij})_{ij}.$$

Obratno lahko vsakemu linearnemu funkcionalu $s \in \mathcal{L}(M_n(\mathcal{M}), \C) $ priredimo preslikavo 
$$\varphi_s \in \mathcal{L} (\mathcal{M}, M_n (\C)),\quad \varphi_s (a) = n \cdot s^{(n)} ([aE_{ij}]_{i, j}).$$
Faktor $\frac{1}{n}$ v definiciji $s_{\varphi}$ smo dodali zato, da je za enotsko preslikavo $\varphi$ tudi $s_{\varphi}$ enotski funkcional.
Preslikavi 
$$\Phi: \mathcal{B} (\mathcal{M}, M_n (\C)) \to (M_n (\mathcal{M}))^*,\quad \varphi \mapsto s_{\varphi}$$
in 
    $$\Psi: (M_n (\mathcal{M}))^* \to \mathcal{B} (\mathcal{M}, M_n (\C)),\quad s \mapsto {\varphi}_s$$
sta linearni in druga drugi inverzni.

    Linearnost je očitna, zato dokažimo, da je $\Psi \circ \Phi = \id_{\mathcal{B} (\mathcal{M}, M_n (\C))}$.
    Vzemimo poljuben $\varphi \in \mathcal{B} (\mathcal{M}, M_n (\C))$. Potem je za vsak $a \in \mathcal{M}$
    \begin{align*}
        \varphi_{s_{\varphi}} (a) &= n \cdot (s_{\varphi}) ^{(n)} ([a E_{ij}]_{i, j}) = n \cdot [s_{\varphi}(a E_{ij})]_{i, j} \\
        &= n \cdot \frac{1}{n}[\varphi (a)_{ij} ]_{i, j} = [\varphi (a)_{ij} ]_{i, j}\\
        &= \varphi(a).
    \end{align*}
    Dokažimo še $\Phi \circ \Psi = \id_{M_n (\mathcal{M})^*}$. Za $A = [a_{ij}]_{i, j} \in M_n (\mathcal{M})$ je 
    \begin{align*}
        s_{\varphi_s} (A) &= \frac{1}{n} \sum_{i, j} ^n \varphi_s (a_{ij})_{ij} = \frac{1}{n} \sum_{i, j} ^n n \cdot (s^{(n)} ([a_{ij} E_{ij}]_{i, j}))_{ij}\\
        &= \sum_{i, j} ^n ([ s(a_{ij} E_{ij})]_{i, j})_{ij} = \sum_{i, j} ^n  s(a_{ij} E_{ij}) = s(A).
    \end{align*}
    Sedaj lahko dokažemo naslednjo trditev.

\begin{trditev}\label{trd:1.7}
    Za linearno preslikavo $\varphi: \mathcal{S} \to M_n (\C)$ so naslednje trditve ekvivalentne:
    \begin{enumerate}
        \item $\varphi$ je povsem pozitivna;
        \item $\varphi$ je $n$-pozitivna;
        \item $s_\varphi$ je pozitiven funkcional.
    \end{enumerate}
\end{trditev}

\begin{dokaz}
    Dokažimo najprej implikacijo $(2.) \Rightarrow (3.)$.
    Naj bo $$e = \begin{bmatrix}
        e_1 \\ \vdots \\ e_n
    \end{bmatrix} \in \C^{n^2}$$
    kot prej. Potem je $$s_{\varphi} (A) = \frac{1}{n} \sum_{i, j} ^n \varphi(a_{ij})_{ij} = \frac{1}{n} e^* \varphi^{(n)}(A) e$$
    za poljubno matriko $A = [a_{ij}]_{i, j} \in M_n (\mathcal{S})$.
    Če je $\varphi^{(n)}$ pozitivna preslikava, potem je $s_\varphi$ pozitiven funkcional.
    
    Sedaj dokažimo še implikacijo $(3.) \Rightarrow (1.)$.
    Če je $s$ pozitiven funkcional, potem je povsem pozitivna preslikava in posledično je tudi 
    $$\varphi_{s} (a) = n \cdot s^{(n)} ([aE_{ij}]_{i, j}) = n \cdot s^{(n)} (e a e^*)$$
    povsem pozitivna. Če je $s_\varphi$ pozitiven funkcional, potem mora biti $\varphi = \varphi_{s_\varphi}$ povsem pozitivna.
\end{dokaz}

\begin{posledica}
    Naj bo $\mathcal{M} \subseteq \mathfrak{A}$ operatorski sistem.
    Za enotsko linearno preslikavo $\varphi: \mathcal{M} \to M_m (\C)$ so naslednje trditve ekvivalentne:
    \begin{enumerate}
        \item $\varphi$ je povsem skrčitev;
        \item $\varphi$ je $n$-skrčitev;
        \item $s_\varphi$ je skrčitev.
    \end{enumerate}
\end{posledica}

\begin{dokaz}
    Implikacija $(2.) \Rightarrow (3.)$ sledi, ker je 
    $$\| s_\varphi (A) \| = \frac{1}{n} \|e^* \varphi^{(n)}(A) e\| \leq \frac{1}{n} \| e\|^2 \| \varphi^{(n)} (A)\| \leq \frac{1}{n} \cdot n \cdot \| A\|  = \| A\|.$$

    Za implikacijo $(3.) \Rightarrow (1.)$ opazimo, da je $s_\varphi$ enotska skrčitev na $M_n (\mathcal{M})$,
    zato jo lahko razširimo do pozitivnega enotskega funkcionala $\widetilde{s_\varphi}: M_n (\mathcal{M}) + M_n (\mathcal{M})^* \to \C$.
    Po trditvi \ref{trd:1.7} je $\varphi_{\widetilde{s_\varphi}}: \mathcal{M} + \mathcal{M}^* \to M_n (\C)$ povsem pozitivna preslikava.
    Dokažimo, da je razširitev $\varphi$.
    Za poljuben $a \in \mathcal{M}$ je 
    $$\varphi_{\widetilde{s_\varphi}} (a) = n \cdot \widetilde{s_{\varphi}} ^{(n)} ([a E_{ij}]_{i, j}) = n \cdot {s_{\varphi}} ^{(n)} ([a E_{ij}]_{i, j}) = \varphi_{s_\varphi} (a) = \varphi(a),$$
    torej je $\varphi_{\widetilde{s_\varphi}}$ v posebnem enotska preslikava.
    Dokazali smo, da je $\varphi_{\widetilde{s_\varphi}}$ enotska povsem pozitivna preslikava, zato je tudi povsem skrčitev.
    Od tod pa sledi, da velja enako za njeno zožitev $\varphi$.
\end{dokaz}

Trditev \ref{trd:1.7} nam omogoča, da podamo trditev \ref{trd:1.6} za preslikave na operatorskem sistemu.

\begin{posledica}
    Naj bo $\mathcal{S} \subseteq M_n (\C)$ operatorski sistem. 
    Za linearno preslikavo $\varphi: \mathcal{S} \to \mathfrak{B}$ sta naslednji
    trditvi ekvivalentni:
    \begin{enumerate}
        \item $\varphi$ je povsem pozitivna;
        \item $\varphi$ je $n$-pozitivna.
    \end{enumerate}
\end{posledica}

\begin{dokaz}
    Denimo, da je $\varphi$ $n$-pozitivna preslikava. Brez škode za splošnost vzemimo $\mathfrak{B} = \bh$.
    Naj bo $[A_{ij}]_{i, j} \in M_m (M_n (\C))$ pozitivna matrika in $x_1, \dots, x_m \in \mathcal{H}$ poljubni vektorji.
    Ponovno definiramo $$V: \C^m \to \mathcal{H},\quad V e_i = x_i$$ in
    $$\psi : \mathcal{S} \to M_m (\C),\quad \psi (A) = V^* \varphi(A) V.$$
    Ker je $\varphi$ $n$-pozitivna, je $\psi$ prav tako $n$-pozitivna in zato po trditvi \ref{trd:1.7} tudi povsem pozitivna.
    Od tod pa sledi 
    \begin{equation*}
        \sum_{i, j} ^m \langle \varphi (A_{ij}) x_j, x_i \rangle = \sum_{i, j} ^m \langle V^* \varphi (A_{ij}) V e_j, e_i \rangle = \sum_{i, j} ^m \langle \psi (A_{ij}) e_j, e_i \rangle \geq 0. \qedhere
    \end{equation*}
\end{dokaz}

Naj bo $\mathcal{S} \subseteq \mathfrak{A}$ poljuben operatorski sistem.
Dokazujemo glavni izrek tega razdelka, ki nam pove, da lahko povsem pozitivno preslikavo $\varphi: \mathcal{S} \to M_n (\C)$
razširimo do povsem pozitivne preslikave $\psi: \mathfrak{A} \to M_n (\C)$.

Najprej to dokažemo za $n = 1$: naslednji rezultat je zgolj različica Hahn-Banachovega izreka za pozitivne funkcionale.

\begin{lema}[Krein-Riesz]\label{lem:1.4}
    Če je $\varphi: \mathcal{S} \to \C$
    pozitiven funkcional, potem ima pozitivno razširitev $\psi: \mathfrak{A} \to \C$, 
    tako da je $\| \varphi \| = \| \varphi \| =  \varphi (1)$.
\end{lema}

\begin{dokaz}
    Ker je $\varphi: \mathcal{S} \to \C$ pozitiven funkcional, velja po lemi \ref{lem:1.1} $\| \varphi \| = \varphi (1)$.
    Po Hahn-Banachovem izreku obstaja razširitev $\varphi$ do $\varphi: \mathfrak{A} \to \C$,
    tako da je $\| \varphi \| = \| \varphi \|$. Od tod sledi 
    $$\| \psi \| = \| \varphi \| = \varphi (1) = \psi(1),$$
    torej je po lemi \ref{lem:1.1} tudi $\psi$ pozitiven funkcional.
\end{dokaz}

\begin{izrek}\label{izr:aux}
    Če je $\varphi: \mathcal{S} \to M_n (\C)$ povsem pozitivna preslikava,
    potem ima povsem pozitivno razširitev $\psi: \mathfrak{A} \to M_n (\C)$,
    tako da velja $\| \psi \|_{\po} = \| \varphi\|_{\po} = \| \varphi (1)\|$.
\end{izrek}

\begin{dokaz}
    Ker je $\varphi$ povsem pozitivna preslikava, je $s_{\varphi}: M_n(\mathcal{S}) \to \C$
    pozitiven funkcional, zato ga lahko po lemi \ref{lem:1.4} razširimo do pozitivnega funkcionala $s: M_n(\mathfrak{A}) \to \C$.
    Potem je $\psi := \varphi_s : \mathfrak{A} \to M_n (\C)$ povsem pozitivna preslikava.

    Dokazati moramo še, da je razširitev $\varphi$. Res, za poljuben $a \in \mathcal{S}$ imamo 
    $$\psi (a) = \varphi_s (a) = n \cdot s^{(n)} ([a E_{ij}]_{i, j}) = n \cdot s_{\varphi}^{(n)} ([a E_{ij}]_{i, j}) = \varphi_{s_{\varphi}} (a) = \varphi(a).$$
    Ker sta $\varphi$ in $\psi$ povsem pozitivna, sledi še 
    \begin{equation*}
        \| \varphi \|_{\po} \leq \| \psi \|_{\po} = \| \psi (1)\| = \|\varphi(1)\| \leq \| \varphi\|_{\po},
    \end{equation*}
    torej $\| \varphi\|_{\po} = \| \psi \|_{\po} = \| \varphi (1)\|$.
\end{dokaz}

\begin{posledica}
    Naj bo $\mathcal{M} \subseteq \mathfrak{A}$ operatorski prostor.
    Če je $\varphi: \mathcal{M} \to M_n (\C)$ enotska povsem skrčitev,
    potem ima povsem pozitivna razširitev $\psi: \mathfrak{A} \to M_n (\C)$.
\end{posledica}

\begin{dokaz}
    Enotsko povsem skrčitev $\varphi: \mathcal{M} \to M_n (\C)$ lahko razširimo do 
    povsem pozitivne preslikave $\tilde{\varphi}: \mathcal{M} + \mathcal{M}^* \to M_n (\C)$,
    ki jo nato lahko razširimo do povsem pozitivne preslikave $\psi: \mathfrak{A} \to M_n (\C)$ po trditvi \ref{trd:1.7}.
\end{dokaz}

Poglavje zaključimo z verzijo izreka \ref{izr:1.2} za operatorske sisteme.

\begin{trditev}
    Naj bo $\mathcal{S} \subseteq M_n (\C)$ operatorski sistem.
    Za linearno preslikavo $\varphi: \mathcal{S} \to M_m (\C)$ so naslednje trditve ekvivalentne:
    \begin{enumerate}
        \item $\varphi$ je povsem pozitivna;
        \item $\varphi$ je $n$-pozitivna;
        \item $\varphi$ je oblike \eqref{eq:1.5}.
    \end{enumerate}
\end{trditev}

\begin{dokaz}
    Dovolj je dokazati implikacijo $(2.) \Rightarrow (3.)$. Če je $\varphi$ $n$-pozitivna,
    je po trditvi \ref{trd:1.7} povsem pozitivna,
    zato jo lahko po trditvi \ref{izr:aux} razširimo do povsem pozitivne preslikave $\psi: M_n (\C) \to M_m (\C)$, ki mora biti oblike \eqref{eq:1.5}.
    Torej je tudi $\varphi$ oblike \eqref{eq:1.5}.
\end{dokaz}

\section{Arvesonov razširitveni izrek}\label{sec:8}

V poglavju \ref{sec:2} smo dokazali, da lahko vsako povsem pozitivno preslikavo
na operatorskem prostoru razširimo do povsem pozitivne preslikave na pripadajočem operatorskem sistemu.
Iz tamkajšnje diskusije je sledilo, da so operatorski sistemi naraven objekt, na katerem so morfizmi povsem pozitivne preslikave.

Naravno vprašanje je, če imamo 
izrek Hahn-Banachovega tipa, ki bi nam povedal, da 
lahko povsem pozitivno preslikavo $$
\varphi : \mathcal{S} \to \bh
$$ na operatorskem sistemu $\mathcal{S} \subseteq \mathfrak{A}$ razširimo do povsem pozitivne preslikave $$
\psi : \mathfrak{A} \to \bh.
$$ na celotni $C^*$-algebri.
To je vsebina Arvesonovega razširitvenega izreka (\cite{arveson1969}).

V končno dimenzionalnem primeru smo to že dokazali v izreku \ref{izr:aux}.
Natančneje, če je $\mathcal{H}$ končno dimenzionalen,
lahko vsako povsem pozitivno preslikavo
$$
\varphi : \mathcal{S} \to \bh
$$
razširimo do povsem pozitivne preslikave na $\mathfrak{A}$.

Neskončno dimenzionalni primer je bistveno zahtevnejši.
Dokaz temelji na konstrukciji ustrezne topologije na prostoru
$\mathcal{B}(\mathcal{S}, \bh)$,
v kateri bo množica
$$
\PP_r(\mathcal{S}, \bh)
:=
\{L \in \mathcal{B}(\mathcal{S}, \bh)
\mid
\text{$L$ je p.p. in $\|L\| \le r$}\}
$$
kompaktna za vsak $r \ge 0$.
Ta topologija bo izhajala iz šibke-$*$ topologije,
zato bo ključna uporaba Banach-Alaoglujevega izreka iz funkcionalne analize (\cite[Theorem V.3.1., str. 130]{conway1990}).

Naš prvi cilj je zato opisati ustrezno preddualno strukturo prostora
$\mathcal{B}(X,Y^*)$.
Naj bosta $X$ in $Y$ Banachova prostora.
Za $x \in X$ in $y \in Y$ definiramo linearen funkcional
$$
x \otimes y : \mathcal{B}(X,Y^*) \to \mathbb{C},
\qquad
(x \otimes y)(L) := L(x)(y).
$$
\begin{lema}
    Za vsaka $x \in X$ in $y \in Y^*$ velja $\| x \otimes y\| = \| x\| \cdot \| y\|$ in posledično
    $x \otimes y \in \mathcal{B}(X, Y^*)^*$.
\end{lema}

\begin{proof}
    Opazimo, da velja $$| x \otimes y (L)| \leq \| L\| \cdot \| x\| \cdot \| y\|,$$
    od koder sledi $\| x \otimes y\| \leq \| x\| \cdot \| y\|$ in $x \otimes y \in \mathcal{B} (X, Y^*)^*$.
    Dokazujemo, da velja tudi $\| x \otimes y \| = \| x\| \cdot \| y\|$. 
    Preko Hahn-Banachovega izreka lahko konstruiramo funkcional 
    $\varphi_x \in X^*$, tako da je $\varphi_x (x) = \| x\|$ in $\| \varphi_x \| = 1$. 
    Podobno lahko definiramo $\varphi_y \in Y^*$, tako da velja $\varphi_y (y) = \| y\|$ in $\| \varphi_y \| = 1$. 
    Nato definiramo preslikavo $$L_{x, y}: X \to Y^*,\quad L_{x, y} (z) = \varphi_x (z) \varphi_y,$$
    za katero velja
    $$\| L_{x, y} (z)\| \leq | \varphi_x (z) | \cdot \| \varphi_y \| \leq \| \varphi_x\| \cdot \| \varphi_ y\| \cdot \| z\| = \| z\|.$$
    To pomeni, da je $\| L_{x, y} \| \leq 1$ in $L_{x, y} \in \mathcal{B}(X, Y^*)$. Po drugi strani pa je $\| L_{x, y} (x)\| = \| x\| \cdot \| \varphi_y \| = \| x\|$
    in zato $\| L_{x, y} \| = 1$.
    Naposled je
    $$x \otimes y (L_{x, y}) = \| x\| \cdot \| y\| = \| x\| \cdot \| y\| \cdot \| L_{x, y} \|,$$
    zato $\| x \otimes y\| = \| x\| \cdot \| y\|$. 
\end{proof}
 
Definiramo prostor
$$Z := \overline{\lin \{x \otimes y\ |\ x \in X, y \in Y\}} \leq \mathcal{B} (X, Y^*)^*.$$
\begin{lema}
    Banachov prostor $\mathcal{B} (X, Y^*)$ je izometrično izomorfen $Z^*$ s preslikavo
    $$\Phi: \mathcal{B}(X, Y^*) \to Z^*,\quad L \mapsto (x \otimes y \mapsto x \otimes y (L)).$$
\end{lema}

\begin{dokaz}
    Najprej dokažimo, da je $\Phi$ izometrija. Vzemimo poljuben $L \in \mathcal{B}(X, Y^*)$.
    Za vsak $x \otimes y \in Z$ velja
    \begin{align*}
        | \Phi (L) (x \otimes y)| &= | x \otimes y (L)|\\
        &= | L(x) (y) |\\
        &\leq \| L(x)\| \cdot \| y\|\\
        &\leq \| L\| \cdot \| x\| \cdot \| y\|\\
        &= \| L \| \cdot \| x \otimes y\|,
    \end{align*}
    od koder sledi $\| \Phi(L) \| \leq \| L\|$. Po drugi strani pa je
    \begin{align*}
        \| L\| &= \sup_{\| x\| = 1} \| L(x) \| \\
        &= \sup_{\| x\| = 1} \sup_{\| y\| = 1} | L(x) (y)|\\
        &= \sup_{\| x\| = 1} \sup_{\| y\| = 1} | x \otimes y (L)|\\
        &= \sup_{\| x\| = 1} \sup_{\| y\| = 1} | \Phi (L) (x \otimes y)|\\
        &\leq \sup_{\| x \otimes y\| = 1} |\Phi (L) (x \otimes y)\| = \| \Phi (L)\|. 
    \end{align*}
    Nazadnje dokažimo še surjektivnost. Fiksirajmo nek $f \in Z^*$.
    Za vse $x \in X$ lahko definiramo funkcional
    $$f_x : Y \to \C,\quad y \mapsto f(x \otimes y).$$
    Ker je $| f_x (y)| \leq \| f\| \cdot \| x\| \cdot \| y\|$, sledi $f_x \in Y^*$.
    Potem je $$L : X \to Y^*,\quad L(x) = f_x.$$
    omejena preslikava in $\Phi(L) = f$. 
\end{dokaz}

Če identificiramo prostor $\mathcal{B} (X, Y^*)$ z $Z^*$, ga lahko opremimo s šibko-$*$ topologijo.
Tej topologiji pravimo \emph{omejeno šibka (OŠ) topologija}.

\begin{lema}
    Naj bo $(L_\lambda)_{\lambda}$ omejena mreža v $\mathcal{B}(X, Y^*)$.
    Potem konvergira $L_\lambda \xrightarrow{} L$ v OŠ topologiji natanko tedaj, ko $L_\lambda (x) \xrightarrow{} L(x)$ v šibki-$*$ topologiji za vse $x \in X$.
\end{lema}

\begin{dokaz}
    Najprej dokažemo implikacijo $(\Rightarrow)$. Če $L_\lambda \to L$ v OŠ topologiji, potem 
    $$L_\lambda (x) (y) = \Phi (L_\lambda) (x \otimes y) \to \Phi (L) (x \otimes y) = L(x) (y)$$
    za vse $x \in X$ in $y \in Y$, torej $L_\lambda (x) \to L (x)$ v šibki-$*$ topologiji na $Y^*$. 
    
    Sedaj dokažimo še implikacijo $(\Leftarrow)$. Ker je mreža $(L_\lambda)_\lambda$ omejena z normo, obstaja $M$, tako da je $\| L_\lambda \| \leq M$ za vsak $\lambda \in \Lambda$.
    Če konvergira $L_\lambda (x) \xrightarrow{} L(x)$ v šibki-$*$ topologiji za vse $x \in X$, tedaj
    $$\Phi (L_\lambda) (x \otimes y) = L_\lambda (x) (y) \to L(x) (y) = \Phi (L) (x \otimes y)$$
    za vse $x \otimes y \in \mathcal{B}(X, Y^*)^*$. Posledično $\Phi (L_\lambda)(z) \to \Phi(L) (z)$ za vse $z \in \lin \{x \otimes y\ |\ x \in X, y \in Y\}$. 
    Po definiciji je $Z = \overline{\{x \otimes y\ |\ x \in X, y \in Y\}}$. Vzemimo poljuben $z \in Z$.
    Potem obstaja $z_0 \in \lin \{x \otimes y\ |\ x \in X, y \in Y\}$, tako da je $\| z - z_0 \| \leq \frac{\varepsilon}{3}$.
    Ker je $\Phi (L_\lambda) (z_0) \to \Phi (L) (z_0)$, obstaja $\lambda_0 \in \Lambda$,
    tako da je $\| \Phi (L_\lambda) (z_0) - \Phi (L) (z_0)\| \leq \frac{\varepsilon}{3M}$ za vse $\lambda \geq \lambda_0$. Potem za vsak $\lambda \geq \lambda_0$ velja tudi 
    \begin{align*}
        \| \Phi (L_\lambda) (z) - \Phi (L) (z)\| &\leq \| \Phi (L_\lambda) (z) - \Phi (L_\lambda) (z_0)\| + \| \Phi (L_\lambda) (z_0) - \Phi (L) (z_0)\|\\
         &+ \|  \Phi (L) (z_0) - \Phi (L) (z)\|\\
        &\leq \| \Phi (L_\lambda) \| \cdot \| z - z_0\| + \frac{\varepsilon}{3} + \| \Phi (L) \| \cdot  \| z_0 - z\|\\
        &\leq \| L_\lambda \| \cdot \| z - z_0\| + \frac{\varepsilon}{3} + \| L \| \cdot  \| z_0 - z\|\\
        &\leq M \cdot \| z - z_0\| + \frac{\varepsilon}{3} + M \cdot  \| z_0 - z\|\\
        &\leq \frac{\varepsilon}{3} + \frac{\varepsilon}{3} + \frac{\varepsilon}{3} = \varepsilon,
    \end{align*}
    torej $\Phi (L_\lambda) (z) \to \Phi (L) (z)$.
\end{dokaz}

Spomnimo, da za poljuben Hilbertov prostor $\mathcal{H}$ obstaja izometrični izomorfizem  
$$\bh \cong L^1 (\mathcal{H})^*,\quad T \mapsto (S \mapsto \tr(TS))$$
pri čemer $L^1 (\mathcal{H})$ označuje Banachov prostor operatorjev s sledjo, opremljen z normo $\| S \|_1 = \tr (|S|)$.

\begin{lema}
    Če je $\{L_\lambda\}_\lambda \subseteq \mathcal{B} (X, \bh)$ omejena mreža,
    potem $L_\lambda \to L$ konvergira v OŠ topologiji natanko tedaj, ko za vse $x \in X$ in $u, v \in \mathcal{H}$ velja
    $\langle L_\lambda (x) u, v \rangle \to \langle L(x) u, v \rangle$.
\end{lema}

\begin{dokaz}
    Za vsaka $u, v \in \mathcal{H}$ definiramo operator ranga $1$ 
    $$S_{u, v} \in \bh,\quad S_{u, v} (\alpha) = \langle \alpha, v \rangle \cdot u.$$
    Opazimo, da velja $\tr (S_{u, v}) = \langle u, v\rangle$.
    Linearne kombinacije operatorjev ranga $1$ tvorijo množico operatorjev končnega ranga,
    ki so gosti v $L^1 (\mathcal{H})$. 

    Sedaj imamo naslednji niz ekvivalenc:
    \begin{align*}
        L_\lambda \to L &\Leftrightarrow L_\lambda (x) \to L (x),\ \forall x \in X \\
        &\Leftrightarrow \tr (L_\lambda (x) S) \to \tr (L (x) S),\ \forall x \in X,\ \forall S \in L^1 (\mathcal{H})\\
        &\Leftrightarrow \tr (L_\lambda (x) S_{u, v}) \to \tr (L (x) S_{u, v}),\ \forall x \in X,\ \forall u, v \in \mathcal{H}\\
        &\Leftrightarrow \langle L_\lambda (x) u, v \rangle \to \langle L (x) u, v \rangle,\ \forall x \in X,\ \forall u, v \in \mathcal{H}. \qedhere
    \end{align*}
\end{dokaz}

\begin{opomba}
    Topologiji na $\bh$, inducirani z družino polnorm 
    $$T \mapsto |\langle Tu, v\rangle|,\quad u, v \in \mathcal{H},$$
    pravimo \emph{šibka operatorska topologija} (od tod naprej ŠOT).
    Prejšnja lema nam torej pove, da omejena mreža $\{L_\lambda\}_{\lambda \in \Lambda} \subseteq \mathcal{B}(X, \bh)$
    konvergira v OŠ topologiji natanko tedaj, ko $\{L_\lambda (x)\}_{\lambda \in \Lambda}$
    konvergira v ŠOT za vsak $x \in X$.
\end{opomba}

\begin{lema}\label{lem:1.5}
    Naj bo $\mathfrak{A}$ $C^*$-algebra in $\mathcal{S} \subseteq \mathfrak{A}$ zaprt operatorski sistem. Potem je 
    $$\PP_r (\mathcal{S}, \bh) := \{L \in \mathcal{B} (\mathcal{S}, \bh)\ |\ \textrm{$L$ p.p.~in$\| L\| \leq r$}\}$$
    kompaktna množica v OŠ topologiji.
\end{lema}

\begin{proof}
    Po Banach-Alaoglujevem izreku je
    $$B_r (\mathcal{S}, \bh) := \{L \in \mathcal{B} (\mathcal{S}, \bh)\ |\ \| L\| \leq r\}$$
    kompaktna množica v OŠ topologiji. Preostane le še pokazati, da je $\PP_r$ zaprt podprostor $B_r$.
    Naj bo $(L_\lambda)_\lambda$ mreža v $\PP_r$, ki konvergira v OŠ topologiji k nekemu $L \in B_r$.
    Dokazati moramo, da je $L$ povsem pozitivna. Naj bo $n \in \N$ in $[a_{ij}]_{i, j} \in M_n (\mathcal{S})$ pozitivna matrika.
    Potem za poljubne $x_1, \dots, x_n \in \mathcal{H}$ velja
    $$\sum_{i, j} ^n \langle L_\lambda  ([a_{ij}]_{i, j})x_j, x_i \rangle \to \sum_{i, j} ^n \langle L  ([a_{ij}]_{i, j}) x_j, x_i \rangle.$$
    Ker je vsota na levi nenegativna za vsak $\lambda$, mora veljati enako za vsoto na desni.
\end{proof}

Sedaj lahko dokažemo glavni izrek poglavja.

\begin{izrek}[Arveson]\label{izr:1.3}
    Naj bo $\mathfrak{A}$ $C^*$-algebra, $\mathcal{S} \subseteq \mathfrak{A}$ operatorski sistem in $\varphi: \mathcal{S} \to \bh$ povsem pozitivna.
    Potem lahko $\varphi$ razširimo do povsem pozitivne preslikave $\psi: \mathfrak{A} \to \bh$, tako da velja
    $\| \psi \|_{\po} = \| \varphi \|_{\po} = \| \varphi (1)\|$.
\end{izrek}

\begin{dokaz}
    Brez škode za splošnost lahko predpostavimo, da je $\mathcal{S}$ zaprt podprostor in zato Banachov.
    Za vsak končno dimenzionalen podprostor $\mathcal{F} \leq \mathcal{H}$
    definiramo povsem pozitivno preslikavo $$\varphi_{\mathcal{F}}: \mathcal{S} \to \mathcal{B}(\mathcal{F}),\quad a \mapsto P_{\mathcal{F}} \varphi(a)\big|_{\mathcal{F}},$$
    kjer je $P_\mathcal{F}: \mathcal{H} \to \mathcal{F}$ projekcija.
    Ker je $\dim (\mathcal{F}) = n < \infty$, lahko identificiramo $\mathcal{B} (\mathcal{F}) \equiv M_n (\C)$.
    Po trditvi \ref{izr:aux} lahko $\varphi_{\mathcal{F}}$ razširimo do povsem pozitivne preslikave $\varphi_\mathcal{F} ' : \mathfrak{A} \to \mathcal{B}(\mathcal{F})$,
    tako da je $\| \varphi_{\mathcal{F}} '\| = \| \varphi_\mathcal{F}\|$.
    Nazadnje za razcep $\mathcal{H} = \mathcal{F} \oplus \mathcal{F}^\perp$ definiramo povsem pozitivno preslikavo
    $$\psi_{\mathcal{F}}: \mathfrak{A} \to \bh,\quad \psi_{\mathcal{F}} (a) = \begin{bmatrix}
        \varphi_{\mathcal{F}} ' (a) & 0\\
        0 & 0
    \end{bmatrix}$$
    in vidimo, da velja $\| \psi_{\mathcal{F}} \| = \| \varphi_{\mathcal{F}} ' \| = \| \varphi_{\mathcal{F}} \| \leq \| \varphi\|$.

    Opazimo, da je $\{\psi_{\mathcal{F}}\}$ mreža v $\PP_{\| \varphi\|} (\mathfrak{A}, \bh)$,
    pri čemer $\mathcal{F}$ teče po vseh končnih podprostorih v $\mathcal{H}$.
    Ker je $\PP_{\| \varphi\|} (\mathfrak{A}, \bh)$ kompaktna množica, ima mreža $\{\psi_{\mathcal{F}}\}$
    stekališče v $\psi \in \PP_{\| \varphi\|} (\mathfrak{A}, \bh)$. 

    Dokazujemo, da je $\psi\big|_{\mathcal{S}} = \varphi$. Fiksirajmo $a \in \mathcal{S}$. 
    Nato vzemimo poljubna vektorja $x, y \in \mathcal{H}$
    in definirajmo prostor $\mathcal{F}_{x, y} := \lin \{x, y\} \leq \mathcal{H}$. Potem za vsak končnorazsežen prostor $\mathcal{K} \subseteq \mathcal{H}$,
    za katerega velja $\mathcal{F} \geq \mathcal{F}_{x, y}$, dobimo 
    $$\langle \psi_{\mathcal{F}}(a) x, y \rangle  = \langle P_{\mathcal{F}} \varphi(a)P_{\mathcal{F}} x, y\rangle = \langle \varphi (a) x, y  \rangle,$$
    torej $\psi_{\mathcal{F}} (a) \to \varphi (a)$ v ŠOT. Po drugi strani pa ima 
    $\{\psi_{\mathcal{F}} (a)\}_{\mathcal{F}}$ stekališče v $\psi(a)$, od koder sledi $\varphi(a) = \psi(a)$.

    Oglejmo si še normo razširitve $\psi$. Ker sta $\varphi$ in $\psi$ povsem pozitivna, velja
    \begin{equation*}
        \| \psi \|_{\po} = \| \psi \| = \| \psi (1) \| = \| \varphi(1) \| = \| \varphi \| = \| \varphi \|_{\po}. \qedhere 
    \end{equation*}
\end{dokaz}

\begin{opomba}
    Predpostavka o povsem pozitivnosti v Arvesonovem razširitvenem izreku je ključna,
    saj v splošnem ne moremo razširiti pozitivne preslikave na operatorskem sistemu do pozitivne preslikave na celotni $C^*$-algebri.
    Vzemimo na primer pozitivno preslikavo $\varphi$ iz zgleda \ref{zgl:1.6}, za katero velja $\| \varphi \| = 2 \cdot \| \varphi(1)\|$.
    Po trditvi \ref{pos:russo-dye} pa za vsako pozitivno preslikavo na $C^*$-algebri velja $\| \varphi \| = \| \varphi(1)\|$.
\end{opomba}

\begin{posledica}
    Naj bo $\mathcal{M} \subseteq \mathfrak{A}$ operatorski prostor in $\varphi: \mathcal{M} \to \bh$ enotska povsem skrčitev.
    Potem lahko $\varphi$ razširimo do enotske povsem pozitivne preslikave $\psi: \mathfrak{A} \to \bh$.
\end{posledica}



